\chapter{Brave New Motivic Spaces}
\label{chap:brave-new-motivic-spaces}

\section{Purpose and standing conventions}
\label{sec:brave-new-motivic-purpose}

This chapter constructs the unstable brave new motivic homotopy theory
associated with the geometric input of
\Ref{chap:smooth-spectral-motivic-site}.  For a quasi-compact
quasi-separated spectral scheme \(S\), the preceding chapter supplied the
essentially small category
\[
    \operatorname{Sm}^{\mathrm{fp}}_{/S},
\]
the spectral Nisnevich topology
\[
    \tau_{\mathrm{Nis}},
\]
and the brave new affine line
\[
    \mathbb A^1_S
    =
    \operatorname{Spec}_S
    \operatorname{Sym}_{\mathcal O_S}(\mathcal O_S).
\]
We impose Nisnevich descent and brave new affine-line invariance by an
accessible localization of the presheaf \(\infty\)-category.  The result
is the presentable \(\infty\)-category
\[
    \mathbf H(S)
\]
of brave new motivic spaces over \(S\).

The chapter then constructs the Cartesian symmetric monoidal structure,
pointed motivic spaces, inverse and direct image for arbitrary morphisms
of quasi-compact quasi-separated spectral schemes, and smooth direct
image for smooth morphisms of finite presentation.  The output is the
unstable
\[
    (*,\sharp,\times)
\]
formalism.  Closed localization, homotopy purity, and nil-invariance are
reserved for the following chapter.

All conventions of the preceding chapters remain in force.  In
particular, sheaves are not hypercompleted unless explicitly stated.
Throughout the chapter,
\[
    S
\]
denotes a quasi-compact quasi-separated spectral scheme and
\[
    \mathcal C_S
    :=
    \operatorname{Sm}^{\mathrm{fp}}_{/S}.
\]
We write
\[
    \operatorname{PSh}(S)
    :=
    \operatorname{PSh}(\mathcal C_S)
    =
    \operatorname{Fun}(\mathcal C_S^{\mathrm{op}},\mathcal S).
\]
The Yoneda embedding is denoted
\[
    h_S
    \colon
    \mathcal C_S
    \hookrightarrow
    \operatorname{PSh}(S).
\]

\section{Presheaves and motivic localization}
\label{sec:brave-new-motivic-localization}

\subsection{The free presentable category on the smooth site}

\begin{proposition}[Presheaves as a free presentable category]
\label{prop:motivic-presheaf-free-cocompletion}
The \(\infty\)-category
\[
    \operatorname{PSh}(S)
\]
is presentable.  Limits and colimits are computed objectwise, and its
Cartesian symmetric monoidal structure is closed and preserves colimits
separately in each variable.

For every presentable \(\infty\)-category \(\mathcal D\), restriction
along the Yoneda embedding induces an equivalence
\[
    \operatorname{Fun}^{L}
    \bigl(
        \operatorname{PSh}(S),
        \mathcal D
    \bigr)
    \xrightarrow{\simeq}
    \operatorname{Fun}
    \bigl(
        \mathcal C_S,
        \mathcal D
    \bigr).
\]
Under this equivalence, a functor
\[
    \Phi
    \colon
    \mathcal C_S
    \longrightarrow
    \mathcal D
\]
corresponds to its colimit-preserving extension
\[
    \overline\Phi
    :=
    \operatorname{Lan}_{h_S}(\Phi).
\]
\end{proposition}

\begin{proof}
The category \(\mathcal C_S\) is essentially small by
\Ref{prop:smooth-site-essentially-small}.  After replacing it by a small
equivalent category, the assertions are the standard properties of a
space-valued presheaf category.  Presentability and objectwise
computation of limits and colimits follow from presentability of
\(\mathcal S\).  Cartesian closedness follows from the existence of
right adjoints to the functors
\[
    F\times -
    \colon
    \operatorname{PSh}(S)
    \longrightarrow
    \operatorname{PSh}(S).
\]
The universal property is the free-cocompletion theorem for the Yoneda
embedding.  See \cite{LurieHTT,LurieHA}.
\end{proof}

\begin{notation}[Internal mapping presheaf]
\label{not:motivic-internal-map}
For presheaves \(F,G\in\operatorname{PSh}(S)\), write
\[
    \underline{\operatorname{Map}}_S(F,G)
\]
for the internal mapping object of the Cartesian closed category
\(\operatorname{PSh}(S)\).  Thus
\[
    \operatorname{Map}_{\operatorname{PSh}(S)}
    \bigl(
        K\times F,
        G
    \bigr)
    \simeq
    \operatorname{Map}_{\operatorname{PSh}(S)}
    \bigl(
        K,
        \underline{\operatorname{Map}}_S(F,G)
    \bigr).
\]
\end{notation}

\subsection{Nisnevich excision morphisms}

\begin{construction}[Nisnevich excision morphisms]
\label{con:nisnevich-excision-morphisms}
Let
\[
\begin{tikzcd}[column sep=large, row sep=large]
    W
        \arrow[r]
        \arrow[d]
    &
    V
        \arrow[d]
    \\
    U
        \arrow[r]
    &
    X
\end{tikzcd}
\]
be a spectral Nisnevich distinguished square in \(\mathcal C_S\).  The
Yoneda image of the square determines a canonical morphism
\[
    e_Q
    \colon
    h_S(U)
    \amalg_{h_S(W)}
    h_S(V)
    \longrightarrow
    h_S(X).
\]
Let
\[
    e_{\varnothing}
    \colon
    \varnothing_{\operatorname{PSh}(S)}
    \longrightarrow
    h_S(\varnothing)
\]
be the canonical morphism from the initial presheaf to the representable
presheaf of the empty spectral scheme.

Write
\[
    E_{\mathrm{Nis},S}
\]
for the set consisting of \(e_{\varnothing}\) and the morphisms \(e_Q\)
as \(Q\) ranges through a set of representatives of spectral Nisnevich
distinguished squares.
\end{construction}

The use of a set is justified by essential smallness of
\(\mathcal C_S\).  Replacing the chosen representatives changes neither
the local objects nor the resulting localization.

\begin{proposition}[Local objects for Nisnevich excision]
\label{prop:nisnevich-local-objects}
For a presheaf \(F\in\operatorname{PSh}(S)\), the following are
equivalent.
\begin{enumerate}[label=(\roman*)]
    \item The presheaf \(F\) is local with respect to every morphism in
    \(E_{\mathrm{Nis},S}\).

    \item The space \(F(\varnothing)\) is contractible and every spectral
    Nisnevich distinguished square is carried to a Cartesian square.

    \item The presheaf \(F\) is a sheaf for
    \(\tau_{\mathrm{Nis}}\).
\end{enumerate}
\end{proposition}

\begin{proof}
Locality with respect to \(e_{\varnothing}\) is the assertion that
\[
    \operatorname{Map}
    \bigl(
        h_S(\varnothing),
        F
    \bigr)
    \longrightarrow
    \operatorname{Map}
    \bigl(
        \varnothing_{\operatorname{PSh}(S)},
        F
    \bigr)
\]
is an equivalence.  By Yoneda, its source is \(F(\varnothing)\), while
its target is contractible.

For a distinguished square \(Q\), mapping \(e_Q\) into \(F\) gives
\[
    F(X)
    \longrightarrow
    F(U)
    \times_{F(W)}
    F(V).
\]
Thus locality with respect to \(e_Q\) is precisely Nisnevich excision
for \(Q\).  This proves the equivalence of (i) and (ii).  The equivalence
of (ii) and (iii) is
\Ref{thm:nisnevich-excision-criterion}.
\end{proof}

\begin{definition}[Nisnevich-local spectral spaces]
\label{def:nisnevich-local-spectral-spaces}
Let
\[
    \operatorname{Spc}_{\mathrm{Nis}}(S)
    \subseteq
    \operatorname{PSh}(S)
\]
be the full subcategory spanned by the
\(E_{\mathrm{Nis},S}\)-local objects.  Equivalently,
\[
    \operatorname{Spc}_{\mathrm{Nis}}(S)
    =
    \operatorname{Sh}_{\tau_{\mathrm{Nis}}}(\mathcal C_S).
\]
\end{definition}

\begin{proposition}[Nisnevich localization]
\label{prop:nisnevich-localization-exists}
The inclusion
\[
    \operatorname{Spc}_{\mathrm{Nis}}(S)
    \hookrightarrow
    \operatorname{PSh}(S)
\]
admits an accessible left exact left adjoint
\[
    L_{\mathrm{Nis},S}
    \colon
    \operatorname{PSh}(S)
    \longrightarrow
    \operatorname{Spc}_{\mathrm{Nis}}(S).
\]
\end{proposition}

\begin{proof}
By \Ref{prop:nisnevich-local-objects}, the local objects are the sheaves
for the Grothendieck topology \(\tau_{\mathrm{Nis}}\).  Sheafification
for a Grothendieck topology on an essentially small \(\infty\)-category
is an accessible left exact localization of the presheaf category.  See
\cite{LurieHTT}.
\end{proof}

\subsection{Affine-line invariance}

\begin{definition}[Brave new affine-line invariance]
\label{def:brave-new-affine-line-invariance}
A presheaf
\[
    F
    \in
    \operatorname{PSh}(S)
\]
is \textbf{\(\mathbb A^1\)-invariant} if, for every
\(X\in\mathcal C_S\), the projection induces an equivalence
\[
    F(X)
    \xrightarrow{\simeq}
    F(X\times_S\mathbb A^1_S).
\]
\end{definition}

Recall from
\Ref{def:affine-line-projection-class}
the set of morphisms
\[
    W_{\mathbb A^1,S}
    =
    \left\{
        h_S(X\times_S\mathbb A^1_S)
        \longrightarrow
        h_S(X)
    \right\}_{X\in\mathcal C_S}.
\]

\begin{proposition}[Affine-line local objects]
\label{prop:affine-line-local-objects}
A presheaf \(F\in\operatorname{PSh}(S)\) is local with respect to
\(W_{\mathbb A^1,S}\) if and only if it is
\(\mathbb A^1\)-invariant.
\end{proposition}

\begin{proof}
For every \(X\in\mathcal C_S\), Yoneda identifies the map obtained by
mapping the projection
\[
    h_S(X\times_S\mathbb A^1_S)
    \longrightarrow
    h_S(X)
\]
into \(F\) with
\[
    F(X)
    \longrightarrow
    F(X\times_S\mathbb A^1_S).
\]
The assertion follows.
\end{proof}

\subsection{The category of brave new motivic spaces}

\begin{definition}[Motivic localization data]
\label{def:brave-new-motivic-localization-data}
Put
\[
    W_{\mathrm{mot},S}
    :=
    E_{\mathrm{Nis},S}
    \cup
    W_{\mathbb A^1,S}.
\]

A presheaf
\[
    F
    \in
    \operatorname{PSh}(S)
\]
is \textbf{motivic-local} if it is local with respect to every
morphism in \(W_{\mathrm{mot},S}\).

A morphism
\[
    u
    \colon
    F
    \longrightarrow
    G
\]
in \(\operatorname{PSh}(S)\) is a
\textbf{motivic equivalence} if, for every motivic-local presheaf
\[
    Z
    \in
    \operatorname{PSh}(S),
\]
the induced morphism
\[
    \operatorname{Map}_{\operatorname{PSh}(S)}
    \bigl(
        G,
        Z
    \bigr)
    \longrightarrow
    \operatorname{Map}_{\operatorname{PSh}(S)}
    \bigl(
        F,
        Z
    \bigr)
\]
is an equivalence.
\end{definition}

\begin{theorem}[Existence and characterization of brave new motivic spaces]
\label{thm:brave-new-motivic-localization}
There is an accessible localization
\[
    L_{\mathrm{mot},S}
    \colon
    \operatorname{PSh}(S)
    \longrightarrow
    \mathbf H(S)
\]
whose local objects are precisely the presheaves \(F\) satisfying:
\begin{enumerate}[label=(\roman*)]
    \item \(F(\varnothing)\) is contractible;

    \item \(F\) carries every spectral Nisnevich distinguished square to
    a Cartesian square;

    \item for every \(X\in\mathcal C_S\), the projection induces an
    equivalence
    \[
        F(X)
        \xrightarrow{\simeq}
        F(X\times_S\mathbb A^1_S).
    \]
\end{enumerate}
Equivalently,
\[
    \mathbf H(S)
    =
    \operatorname{Sh}_{\mathrm{Nis}}(\mathcal C_S)
    \cap
    \operatorname{PSh}(S)^{\mathbb A^1}
\]
as a full subcategory of \(\operatorname{PSh}(S)\).

The \(\infty\)-category \(\mathbf H(S)\) is presentable, and the
inclusion
\[
    \iota_S
    \colon
    \mathbf H(S)
    \hookrightarrow
    \operatorname{PSh}(S)
\]
is fully faithful and accessible.

Moreover, the class of motivic equivalences in
\(\operatorname{PSh}(S)\) is precisely the strongly saturated class
generated by
\[
    W_{\mathrm{mot},S}.
\]
\end{theorem}

\begin{proof}
The set
\[
    W_{\mathrm{mot},S}
    =
    E_{\mathrm{Nis},S}
    \cup
    W_{\mathbb A^1,S}
\]
is a set of morphisms in the presentable category
\(\operatorname{PSh}(S)\). The accessible localization theorem
therefore supplies an accessible reflective localization
\[
    L_{\mathrm{mot},S}
    \colon
    \operatorname{PSh}(S)
    \longrightarrow
    \mathbf H(S)
\]
whose essential image is the full subcategory of
\(W_{\mathrm{mot},S}\)-local objects. In particular,
\(\mathbf H(S)\) is presentable, and its inclusion
\[
    \iota_S
    \colon
    \mathbf H(S)
    \hookrightarrow
    \operatorname{PSh}(S)
\]
is fully faithful and accessible.

By
\Ref{prop:nisnevich-local-objects},
locality with respect to
\[
    E_{\mathrm{Nis},S}
\]
is equivalent to the conditions that \(F(\varnothing)\) be
contractible and that every spectral Nisnevich distinguished square be
carried to a Cartesian square.

By
\Ref{prop:affine-line-local-objects},
locality with respect to
\[
    W_{\mathbb A^1,S}
\]
is equivalent to brave new affine-line invariance. Thus the
\(W_{\mathrm{mot},S}\)-local objects are precisely the presheaves
satisfying conditions (i)--(iii).

Equivalently, the full subcategory of local objects is
\[
    \operatorname{Sh}_{\mathrm{Nis}}(\mathcal C_S)
    \cap
    \operatorname{PSh}(S)^{\mathbb A^1}.
\]

Finally, for an accessible localization of a presentable
\(\infty\)-category generated by a set of morphisms, the class of
morphisms inverted by the localization is the strongly saturated class
generated by that set. Consequently, the motivic equivalences are
precisely the strongly saturated class generated by
\[
    W_{\mathrm{mot},S}.
\]

See
\cite{LurieHTT,KhanBraveNewMotivicI}.
\end{proof}

\begin{definition}[Brave new motivic space]
\label{def:brave-new-motivic-space}
An object of
\[
    \mathbf H(S)
\]
is called a \textbf{brave new motivic space over \(S\)}.
\end{definition}

\begin{proposition}[Two-stage motivic localization]
\label{prop:two-stage-brave-new-motivic-localization}
Let
\[
    L_{\mathbb A^1,S}
    \colon
    \operatorname{Spc}_{\mathrm{Nis}}(S)
    \longrightarrow
    \operatorname{Spc}_{\mathrm{Nis}}(S)^{\mathbb A^1}
\]
be the accessible localization at the images under
\(L_{\mathrm{Nis},S}\) of the morphisms in
\(W_{\mathbb A^1,S}\).  Then there are canonical equivalences
\[
    \mathbf H(S)
    \simeq
    \operatorname{Spc}_{\mathrm{Nis}}(S)^{\mathbb A^1}
\]
and
\[
    L_{\mathrm{mot},S}
    \simeq
    L_{\mathbb A^1,S}
    L_{\mathrm{Nis},S}.
\]
\end{proposition}

\begin{proof}
A local object for the second localization is a Nisnevich-local object
which is local with respect to the images of
\(W_{\mathbb A^1,S}\).  By adjunction, this is equivalent to being local
with respect to \(W_{\mathbb A^1,S}\) as an object of
\(\operatorname{PSh}(S)\).  Hence the local objects for the composite
are exactly the \(W_{\mathrm{mot},S}\)-local objects.  Reflective
localizations with the same local objects are canonically equivalent.
\end{proof}

\begin{definition}[Motive of a smooth spectral scheme]
\label{def:brave-new-smooth-motive}
For
\[
    X
    \in
    \mathcal C_S,
\]
define its \textbf{unstable brave new motive} by
\[
    \mathrm M_S(X)
    :=
    L_{\mathrm{mot},S}h_S(X)
    \in
    \mathbf H(S).
\]
The resulting functor is denoted
\[
    \mathrm M_S
    \colon
    \mathcal C_S
    \longrightarrow
    \mathbf H(S).
\]
\end{definition}

Since the spectral Nisnevich topology is subcanonical,
\Ref{thm:spectral-nisnevich-subcanonical}
shows that
\[
    L_{\mathrm{Nis},S}h_S(X)
    \simeq
    h_S(X).
\]
Thus \(\mathrm M_S(X)\) is obtained from the representable presheaf by
imposing only the remaining \(\mathbb A^1\)-local equivalences inside
the Nisnevich sheaf category.

\begin{proposition}[Evaluation represented by smooth motives]
\label{prop:motivic-evaluation-represented}
For every \(X\in\mathcal C_S\) and every
\(F\in\mathbf H(S)\), there is a natural equivalence
\[
    \operatorname{Map}_{\mathbf H(S)}
    \bigl(
        \mathrm M_S(X),
        F
    \bigr)
    \simeq
    F(X).
\]
\end{proposition}

\begin{proof}
The localization adjunction and Yoneda give
\[
\begin{aligned}
    \operatorname{Map}_{\mathbf H(S)}
    \bigl(
        L_{\mathrm{mot},S}h_S(X),
        F
    \bigr)
    &\simeq
    \operatorname{Map}_{\operatorname{PSh}(S)}
    \bigl(
        h_S(X),
        \iota_SF
    \bigr)
    \\
    &\simeq
    F(X).
\end{aligned}
\]
\end{proof}

\begin{theorem}[Universal property of brave new motivic spaces]
\label{thm:brave-new-motivic-universal-property}
Let \(\mathcal D\) be a presentable \(\infty\)-category.  Restriction
along
\[
    \mathrm M_S
    \colon
    \mathcal C_S
    \longrightarrow
    \mathbf H(S)
\]
induces an equivalence between:
\begin{enumerate}[label=(\roman*)]
    \item colimit-preserving functors
    \[
        \Psi
        \colon
        \mathbf H(S)
        \longrightarrow
        \mathcal D;
    \]

    \item functors
    \[
        \Phi
        \colon
        \mathcal C_S
        \longrightarrow
        \mathcal D
    \]
    satisfying the following conditions:
    \begin{enumerate}[label=(\alph*)]
        \item \(\Phi(\varnothing)\) is initial in \(\mathcal D\);

        \item for every spectral Nisnevich distinguished square
        \[
        \begin{tikzcd}[column sep=large, row sep=large]
            W \arrow[r] \arrow[d]
            &
            V \arrow[d]
            \\
            U \arrow[r]
            &
            X,
        \end{tikzcd}
        \]
        the induced square
        \[
        \begin{tikzcd}[column sep=large, row sep=large]
            \Phi(W) \arrow[r] \arrow[d]
            &
            \Phi(V) \arrow[d]
            \\
            \Phi(U) \arrow[r]
            &
            \Phi(X)
        \end{tikzcd}
        \]
        is cocartesian;

        \item for every \(X\in\mathcal C_S\), the morphism
        \[
            \Phi(X\times_S\mathbb A^1_S)
            \longrightarrow
            \Phi(X)
        \]
        is an equivalence.
    \end{enumerate}
\end{enumerate}
\end{theorem}

\begin{proof}
By
\Ref{prop:motivic-presheaf-free-cocompletion},
a functor
\[
    \Phi
    \colon
    \mathcal C_S
    \longrightarrow
    \mathcal D
\]
extends uniquely to a colimit-preserving functor
\[
    \overline\Phi
    \colon
    \operatorname{PSh}(S)
    \longrightarrow
    \mathcal D.
\]
The extension factors through the localization
\[
    L_{\mathrm{mot},S}
\]
if and only if it carries every morphism in
\(W_{\mathrm{mot},S}\) to an equivalence.

The image of
\[
    e_{\varnothing}
    \colon
    \varnothing_{\operatorname{PSh}(S)}
    \longrightarrow
    h_S(\varnothing)
\]
is
\[
    0_{\mathcal D}
    \longrightarrow
    \Phi(\varnothing),
\]
so it is invertible exactly when (a) holds.

Since \(\overline\Phi\) preserves colimits, the image of a Nisnevich
excision morphism is
\[
    \Phi(U)
    \amalg_{\Phi(W)}
    \Phi(V)
    \longrightarrow
    \Phi(X).
\]
It is invertible exactly when the displayed square in (b) is
cocartesian.

Finally, the images of the morphisms in
\(W_{\mathbb A^1,S}\) are exactly the morphisms in (c).
The localization universal property now proves the assertion.
\end{proof}

\begin{corollary}[Finite coproducts under the motivic Yoneda functor]
\label{cor:motivic-yoneda-finite-coproducts}
The functor
\[
    \mathrm M_S
    \colon
    \mathcal C_S
    \longrightarrow
    \mathbf H(S)
\]
preserves finite coproducts.
\end{corollary}

\begin{proof}
The empty spectral scheme is carried to an initial object because
\[
    \operatorname{Map}_{\mathbf H(S)}
    \bigl(
        \mathrm M_S(\varnothing),
        F
    \bigr)
    \simeq
    F(\varnothing)
    \simeq
    *
\]
for every \(F\in\mathbf H(S)\).

Let \(X,Y\in\mathcal C_S\).  The square
\[
\begin{tikzcd}[column sep=large, row sep=large]
    \varnothing
        \arrow[r]
        \arrow[d]
    &
    Y
        \arrow[d]
    \\
    X
        \arrow[r]
    &
    X\amalg Y
\end{tikzcd}
\]
is a spectral Nisnevich distinguished square: the two component
inclusions are open immersions, and \(Y\) identifies with the closed
complement of \(X\).  The universal property of
\Ref{thm:brave-new-motivic-universal-property}
therefore gives
\[
    \mathrm M_S(X)
    \amalg
    \mathrm M_S(Y)
    \xrightarrow{\simeq}
    \mathrm M_S(X\amalg Y).
\]
\end{proof}

\begin{proposition}[Compact generation]
\label{prop:brave-new-motivic-compact-generation}
The \(\infty\)-category \(\mathbf H(S)\) is compactly generated by the
set
\[
    \left\{
        \mathrm M_S(X)
    \right\}_{X\in\mathcal C_S}.
\]
In particular, every \(\mathrm M_S(X)\) is compact.
\end{proposition}

\begin{proof}
Representable presheaves are compact in
\(\operatorname{PSh}(S)\), since evaluation at any
\(X\in\mathcal C_S\) preserves all colimits.

The motivic-local conditions of
\Ref{thm:brave-new-motivic-localization}
are preserved under filtered colimits. Indeed, filtered colimits of
spaces commute with finite limits, so the conditions
\[
    F(\varnothing)
    \simeq
    *
\]
and
\[
    F(X)
    \xrightarrow{\simeq}
    F(U)
    \times_{F(W)}
    F(V)
\]
for spectral Nisnevich distinguished squares are preserved under
filtered colimits. Brave new affine-line invariance is also preserved
under filtered colimits because it is an objectwise equivalence
condition.

Consequently, the inclusion
\[
    \iota_S
    \colon
    \mathbf H(S)
    \hookrightarrow
    \operatorname{PSh}(S)
\]
preserves filtered colimits.

Let
\[
    \{F_\alpha\}_{\alpha\in I}
\]
be a filtered diagram in \(\mathbf H(S)\). Since its colimit is already
motivic-local, its colimit in \(\mathbf H(S)\) is computed by its
objectwise colimit in \(\operatorname{PSh}(S)\). Therefore
\[
\begin{aligned}
    \operatorname{Map}_{\mathbf H(S)}
    \left(
        \mathrm M_S(X),
        \operatorname*{colim}_{\alpha}F_\alpha
    \right)
    &\simeq
    \left(
        \operatorname*{colim}_{\alpha}F_\alpha
    \right)(X)
    \\
    &\simeq
    \operatorname*{colim}_{\alpha}F_\alpha(X)
    \\
    &\simeq
    \operatorname*{colim}_{\alpha}
    \operatorname{Map}_{\mathbf H(S)}
    \bigl(
        \mathrm M_S(X),
        F_\alpha
    \bigr).
\end{aligned}
\]
Thus every \(\mathrm M_S(X)\) is compact.

The family
\[
    \left\{
        \mathrm M_S(X)
    \right\}_{X\in\mathcal C_S}
\]
detects equivalences. Indeed, let
\[
    \varphi
    \colon
    F
    \longrightarrow
    G
\]
be a morphism in \(\mathbf H(S)\). By
\Ref{prop:motivic-evaluation-represented},
the induced morphism
\[
    \operatorname{Map}_{\mathbf H(S)}
    \bigl(
        \mathrm M_S(X),
        F
    \bigr)
    \longrightarrow
    \operatorname{Map}_{\mathbf H(S)}
    \bigl(
        \mathrm M_S(X),
        G
    \bigr)
\]
identifies with
\[
    F(X)
    \longrightarrow
    G(X).
\]
If this is an equivalence for every \(X\in\mathcal C_S\), then
\(\varphi\) is an objectwise equivalence of presheaves and hence an
equivalence in \(\mathbf H(S)\).

The indexing category \(\mathcal C_S\) is essentially small, so after
choosing a small equivalent category the displayed family is a set. A
presentable \(\infty\)-category admitting a set of compact objects which
detects equivalences is compactly generated by those objects.
\end{proof}

\section{Cartesian structure on brave new motivic spaces}
\label{sec:brave-new-motivic-cartesian-structure}

\subsection{Product stability of motivic equivalences}

\begin{proposition}[Product stability of the generating morphisms]
\label{prop:motivic-generators-product-stable}
Let
\[
    w
    \in
    W_{\mathrm{mot},S}
\]
and let \(Y\in\mathcal C_S\). Then
\[
    w\times h_S(Y)
\]
is a motivic equivalence.
\end{proposition}

\begin{proof}
Suppose first that
\[
    w
    =
    e_Q
\]
is associated with a chosen representative
\[
\begin{tikzcd}[column sep=large, row sep=large]
    W
        \arrow[r]
        \arrow[d]
    &
    V
        \arrow[d]
    \\
    U
        \arrow[r]
    &
    X
\end{tikzcd}
\]
of a spectral Nisnevich distinguished square.

Cartesian product with \(h_S(Y)\) preserves colimits, and the Yoneda
embedding preserves finite products. Hence
\[
    e_Q\times h_S(Y)
\]
identifies with the excision morphism
\[
    h_S(U\times_SY)
    \amalg_{h_S(W\times_SY)}
    h_S(V\times_SY)
    \longrightarrow
    h_S(X\times_SY)
\]
associated with the base-changed square
\[
\begin{tikzcd}[column sep=large, row sep=large]
    W\times_SY
        \arrow[r]
        \arrow[d]
    &
    V\times_SY
        \arrow[d]
    \\
    U\times_SY
        \arrow[r]
    &
    X\times_SY.
\end{tikzcd}
\]
This is again a spectral Nisnevich distinguished square by
\Ref{prop:nisnevich-square-base-change}.

Its excision morphism is equivalent in the arrow category of
\(\operatorname{PSh}(S)\) to the excision morphism attached to the
chosen representative of its equivalence class. The latter belongs to
\(E_{\mathrm{Nis},S}\). Therefore
\[
    e_Q\times h_S(Y)
\]
is a motivic equivalence.

For
\[
    w
    =
    e_{\varnothing},
\]
one has
\[
    \varnothing_{\operatorname{PSh}(S)}
    \times
    h_S(Y)
    \simeq
    \varnothing_{\operatorname{PSh}(S)}
\]
and
\[
    h_S(\varnothing)\times h_S(Y)
    \simeq
    h_S(\varnothing\times_SY)
    \simeq
    h_S(\varnothing).
\]
Thus
\[
    e_{\varnothing}\times h_S(Y)
\]
identifies with \(e_{\varnothing}\), and hence is a motivic
equivalence.

Finally, suppose that \(w\) is the affine-line projection associated
with \(X\):
\[
    h_S(X\times_S\mathbb A^1_S)
    \longrightarrow
    h_S(X).
\]
Its product with \(h_S(Y)\) identifies with
\[
    h_S
    \bigl(
        (X\times_SY)\times_S\mathbb A^1_S
    \bigr)
    \longrightarrow
    h_S(X\times_SY),
\]
which belongs to \(W_{\mathbb A^1,S}\). It is therefore a motivic
equivalence.
\end{proof}

\begin{theorem}[Cartesian monoidal motivic localization]
\label{thm:cartesian-monoidal-motivic-localization}
Motivic equivalences are stable under Cartesian product with arbitrary
presheaves.  The localization
\[
    L_{\mathrm{mot},S}
    \colon
    \operatorname{PSh}(S)
    \longrightarrow
    \mathbf H(S)
\]
is a Cartesian symmetric monoidal localization.  In particular:
\begin{enumerate}[label=(\roman*)]
    \item the Cartesian product of motivic-local presheaves is
    motivic-local;

    \item \(L_{\mathrm{mot},S}\) preserves finite products;

    \item \(\mathbf H(S)\) is Cartesian closed;

    \item the Cartesian product on \(\mathbf H(S)\) preserves colimits
    separately in each variable.
\end{enumerate}
\end{theorem}

\begin{proof}
Fix a presheaf \(K\).  The class of morphisms \(w\) such that
\[
    w\times K
\]
is a motivic equivalence is strongly saturated: Cartesian product with
\(K\) preserves colimits, composition, and retracts.  It therefore
suffices to prove the assertion when \(w\) belongs to
\(W_{\mathrm{mot},S}\).

Write \(K\) as a colimit of representable presheaves.  Cartesian product
in spaces preserves colimits separately, hence
\[
    w\times K
\]
is the corresponding colimit of the morphisms
\[
    w\times h_S(Y).
\]
Each of these is a motivic equivalence by
\Ref{prop:motivic-generators-product-stable}.  The class of local
equivalences is closed under colimits in the arrow category, since
mapping a colimit of arrows into a local object gives a limit of
equivalences.  Thus \(w\times K\) is a motivic equivalence.

The terminal presheaf is \(h_S(S)\).  It is Nisnevich-local and
\(\mathbb A^1\)-invariant, hence motivic-local.  If \(F\) and \(G\) are
motivic-local, then for every motivic equivalence \(A\to B\),
\[
\begin{aligned}
    \operatorname{Map}(B,F\times G)
    &\simeq
    \operatorname{Map}(B,F)
    \times
    \operatorname{Map}(B,G)
    \\
    &\xrightarrow{\simeq}
    \operatorname{Map}(A,F)
    \times
    \operatorname{Map}(A,G)
    \\
    &\simeq
    \operatorname{Map}(A,F\times G).
\end{aligned}
\]
Thus \(F\times G\) is local.

The morphism
\[
    F\times G
    \longrightarrow
    L_{\mathrm{mot},S}F
    \times
    L_{\mathrm{mot},S}G
\]
is a motivic equivalence by product stability, and its target is local.
It therefore exhibits the target as the localization of \(F\times G\).
This proves preservation of finite products.

Let \(F\) be local and let \(K\) be any presheaf.  For every motivic
equivalence \(A\to B\), there are natural equivalences
\[
\begin{aligned}
    \operatorname{Map}
    \bigl(
        B,
        \underline{\operatorname{Map}}_S(K,F)
    \bigr)
    &\simeq
    \operatorname{Map}(B\times K,F)
    \\
    &\xrightarrow{\simeq}
    \operatorname{Map}(A\times K,F)
    \\
    &\simeq
    \operatorname{Map}
    \bigl(
        A,
        \underline{\operatorname{Map}}_S(K,F)
    \bigr).
\end{aligned}
\]
Hence the presheaf internal mapping object is local.  It therefore
restricts to an internal mapping object in \(\mathbf H(S)\), proving
Cartesian closedness.  In a Cartesian closed presentable category,
product with a fixed object is a left adjoint, so it preserves colimits.
See \cite{LurieHA,KhanBraveNewMotivicI}.
\end{proof}

\begin{notation}[Cartesian product and internal maps]
\label{not:motivic-cartesian-product}
We denote the Cartesian product in \(\mathbf H(S)\) by
\[
    F\times_SG
\]
when the base needs to be displayed, and simply by \(F\times G\)
otherwise.  Its unit is the terminal motivic space
\[
    \mathbf 1_S
    :=
    \mathrm M_S(S)
    \simeq
    h_S(S).
\]
The internal mapping object is denoted
\[
    \underline{\operatorname{Map}}_{\mathbf H(S)}(F,G).
\]
\end{notation}

\begin{corollary}[Products of smooth motives]
\label{cor:products-of-smooth-motives}
For \(X,Y\in\mathcal C_S\), there is a canonical equivalence
\[
    \mathrm M_S(X)
    \times_S
    \mathrm M_S(Y)
    \simeq
    \mathrm M_S(X\times_SY).
\]
\end{corollary}

\begin{proof}
The Yoneda embedding preserves products, so
\[
    h_S(X)\times h_S(Y)
    \simeq
    h_S(X\times_SY).
\]
Apply the finite-product-preserving localization
\(L_{\mathrm{mot},S}\).
\end{proof}

\begin{proposition}[Tensoring by spaces]
\label{prop:motivic-tensored-over-spaces}
The presentable category \(\mathbf H(S)\) is canonically tensored and
cotensored over spaces.  For \(K\in\mathcal S\) and
\(F,G\in\mathbf H(S)\), there are natural equivalences
\[
    \operatorname{Map}_{\mathbf H(S)}
    \bigl(
        K\otimes F,
        G
    \bigr)
    \simeq
    \operatorname{Map}_{\mathcal S}
    \left(
        K,
        \operatorname{Map}_{\mathbf H(S)}(F,G)
    \right).
\]
The tensor is compatible with Cartesian products:
\[
    K\otimes(F\times G)
    \simeq
    (K\otimes F)\times G
\]
after interpreting the left tensor through colimits in the first
variable.
\end{proposition}

\begin{proof}
Every presentable \(\infty\)-category is canonically tensored and
cotensored over spaces.  The compatibility follows because Cartesian
product in \(\mathbf H(S)\) preserves colimits separately.
\end{proof}

\section{Pointed brave new motivic spaces}
\label{sec:pointed-brave-new-motivic-spaces}

\subsection{Pointed objects}

\begin{definition}[Pointed brave new motivic spaces]
\label{def:pointed-brave-new-motivic-spaces}
Define
\[
    \mathbf H_\bullet(S)
    :=
    \mathbf H(S)_{*/}
\]
to be the \(\infty\)-category of pointed objects of
\(\mathbf H(S)\).  An object is a pair
\[
    (F,s),
    \qquad
    s
    \colon
    *
    \longrightarrow
    F.
\]
\end{definition}

\begin{proposition}[Presentability and free pointing]
\label{prop:pointed-motivic-presentability}
The category \(\mathbf H_\bullet(S)\) is presentable.  The forgetful
functor
\[
    U_S
    \colon
    \mathbf H_\bullet(S)
    \longrightarrow
    \mathbf H(S)
\]
admits a left adjoint
\[
    (-)_+
    \colon
    \mathbf H(S)
    \longrightarrow
    \mathbf H_\bullet(S),
    \qquad
    F
    \longmapsto
    F_+:=F\amalg *.
\]
The basepoint of \(F_+\) is the second coproduct inclusion.
\end{proposition}

\begin{proof}
An undercategory of a presentable \(\infty\)-category is presentable.
For a pointed object \((G,s)\), the universal property of the coproduct
gives
\[
    \operatorname{Map}_{\mathbf H_\bullet(S)}
    \bigl(
        F_+,
        (G,s)
    \bigr)
    \simeq
    \operatorname{Map}_{\mathbf H(S)}(F,G).
\]
Thus \((-)_{+}\dashv U_S\).
\end{proof}

\begin{definition}[Pointed smooth motive]
\label{def:pointed-smooth-motive}
For \(X\in\mathcal C_S\), put
\[
    \mathrm M_{S,+}(X)
    :=
    \mathrm M_S(X)_+.
\]
\end{definition}

\begin{proposition}[Generation of pointed motivic spaces]
\label{prop:pointed-motivic-generation}
The category \(\mathbf H_\bullet(S)\) is generated under colimits by the
objects
\[
    \mathrm M_{S,+}(X),
    \qquad
    X\in\mathcal C_S.
\]
\end{proposition}

\begin{proof}
Consider the free-pointing adjunction
\[
    (-)_+
    \colon
    \mathbf H(S)
    \rightleftarrows
    \mathbf H_\bullet(S)
    :
    U_S.
\]
The forgetful functor \(U_S\) is monadic, with associated monad
\[
    T
    \colon
    \mathbf H(S)
    \longrightarrow
    \mathbf H(S),
    \qquad
    T(F)
    =
    F\amalg *.
\]

By the monadic bar construction, every pointed motivic space
\[
    G
    \in
    \mathbf H_\bullet(S)
\]
is the geometric realization of a simplicial object whose terms are
free pointed objects. Explicitly, its terms have the form
\[
    \bigl(
        T^nU_SG
    \bigr)_+,
    \qquad
    n\geq0.
\]
Consequently, every object of \(\mathbf H_\bullet(S)\) belongs to the
full subcategory generated under colimits by the free pointed objects
\[
    F_+,
    \qquad
    F\in\mathbf H(S).
\]

By
\Ref{prop:brave-new-motivic-compact-generation},
the category \(\mathbf H(S)\) is generated under colimits by the
objects
\[
    \mathrm M_S(X),
    \qquad
    X\in\mathcal C_S.
\]
Since the free-pointing functor
\[
    (-)_+
\]
is a left adjoint, it preserves colimits. Therefore every free pointed
object \(F_+\) is generated under colimits by
\[
    \mathrm M_S(X)_+
    =
    \mathrm M_{S,+}(X).
\]
It follows that the objects
\[
    \mathrm M_{S,+}(X),
    \qquad
    X\in\mathcal C_S,
\]
generate \(\mathbf H_\bullet(S)\) under colimits.
\end{proof}

\subsection{Smash products}

\begin{construction}[Wedge and smash products]
\label{con:pointed-motivic-smash-product}
Let \(F,G\in\mathbf H_\bullet(S)\).  Their wedge is the coproduct in the
pointed category:
\[
    F\vee G
    :=
    F\amalg_*G.
\]
Their smash product is the pushout
\[
    F\wedge_SG
    :=
    (F\times_SG)
    \amalg_{F\vee G}
    *.
\]
Equivalently,
\[
    F\wedge_SG
    =
    \frac{F\times_SG}{F\vee G}.
\]
\end{construction}

\begin{theorem}[Presentable symmetric monoidal structure on pointed motives]
\label{thm:pointed-motivic-smash-monoidal}
The smash product equips
\[
    \mathbf H_\bullet(S)
\]
with a closed presentably symmetric monoidal structure.  Its tensor unit
is
\[
    \mathbb S^0_S
    :=
    (\mathbf 1_S)_+
    =
    *\amalg *.
\]
The smash product preserves colimits separately in each variable, and
there are natural equivalences
\[
    F_+\wedge_SG_+
    \simeq
    (F\times_SG)_+.
\]
In particular,
\[
    \mathrm M_{S,+}(X)
    \wedge_S
    \mathrm M_{S,+}(Y)
    \simeq
    \mathrm M_{S,+}(X\times_SY).
\]
\end{theorem}

\begin{proof}
The Cartesian product on \(\mathbf H(S)\) preserves colimits separately
by
\Ref{thm:cartesian-monoidal-motivic-localization}.
The standard pointed-object construction therefore equips
\(\mathbf H(S)_{*/}\) with its smash-product symmetric monoidal
structure.  Closedness and separate preservation of colimits follow
from presentability.  The identity
\[
    F_+\wedge G_+
    \simeq
    (F\times G)_+
\]
follows directly from the pushout defining the smash product.  Apply
\Ref{cor:products-of-smooth-motives}
for the final assertion.  See \cite{LurieHA}.
\end{proof}

\begin{notation}[Pointed internal maps]
\label{not:pointed-motivic-internal-map}
The internal mapping object for the smash product is denoted
\[
    \underline{\operatorname{Map}}_{\mathbf H_\bullet(S)}(F,G).
\]
Thus
\[
    \operatorname{Map}_{\mathbf H_\bullet(S)}
    \bigl(
        K\wedge_SF,
        G
    \bigr)
    \simeq
    \operatorname{Map}_{\mathbf H_\bullet(S)}
    \bigl(
        K,
        \underline{\operatorname{Map}}_{\mathbf H_\bullet(S)}(F,G)
    \bigr).
\]
\end{notation}

\subsection{Cofibres, suspension, and loops}

\begin{definition}[Motivic quotient]
\label{def:motivic-quotient}
Let
\[
    u
    \colon
    F
    \longrightarrow
    G
\]
be a morphism in \(\mathbf H(S)\).  Its pointed quotient is
\[
    G/F
    :=
    G
    \amalg_F
    *
    \in
    \mathbf H_\bullet(S).
\]
For a morphism
\[
    U
    \longrightarrow
    X
\]
in \(\mathcal C_S\), write
\[
    X/U
    :=
    \mathrm M_S(X)/
    \mathrm M_S(U).
\]
\end{definition}

\begin{definition}[Suspension and loop objects]
\label{def:motivic-suspension-loop}
For \(F\in\mathbf H_\bullet(S)\), define
\[
    \Sigma F
    :=
    *
    \amalg_F
    *,
\]
and define
\[
    \Omega F
    :=
    *
    \times_F
    *.
\]
\end{definition}

\begin{proposition}[Suspension-loop adjunction]
\label{prop:motivic-suspension-loop-adjunction}
There is a natural adjunction
\[
    \Sigma
    \colon
    \mathbf H_\bullet(S)
    \rightleftarrows
    \mathbf H_\bullet(S)
    :
    \Omega.
\]
Equivalently,
\[
    \Sigma F
    \simeq
    S^1\otimes F
\]
for the canonical tensoring of
\(\mathbf H_\bullet(S)\) over pointed spaces.
\end{proposition}

\begin{proof}
The assertion is the suspension-loop adjunction in any pointed
presentable \(\infty\)-category.  The pushout and pullback formulas are
the colimit and limit descriptions of suspension and loops.
\end{proof}

\section{Functoriality in the spectral base}
\label{sec:brave-new-motivic-base-functoriality}

For the remainder of this section, let
\[
    f
    \colon
    T
    \longrightarrow
    S
\]
be a morphism of quasi-compact quasi-separated spectral schemes.  The
geometric base-change functor of
\Ref{prop:smooth-site-geometric-base-change}
is
\[
    u_f
    :=
    f^*_{\mathrm{geom}}
    \colon
    \mathcal C_S
    \longrightarrow
    \mathcal C_T,
    \qquad
    X
    \longmapsto
    X_T:=X\times_ST.
\]

\subsection{Presheaf inverse and direct image}

\begin{construction}[Presheaf inverse and direct image]
\label{con:presheaf-base-change-adjunction}
Define
\[
    f^*_{\operatorname{PSh}}
    :=
    \operatorname{Lan}_{u_f^{\mathrm{op}}}
    \colon
    \operatorname{PSh}(S)
    \longrightarrow
    \operatorname{PSh}(T)
\]
and
\[
    f_{*,\operatorname{PSh}}
    :=
    (-)\circ u_f^{\mathrm{op}}
    \colon
    \operatorname{PSh}(T)
    \longrightarrow
    \operatorname{PSh}(S).
\]
Then
\[
    f^*_{\operatorname{PSh}}
    \dashv
    f_{*,\operatorname{PSh}}.
\]
\end{construction}

\begin{proposition}[Effect on representables and sections]
\label{prop:presheaf-base-change-formulas}
For every \(X\in\mathcal C_S\), there is a canonical equivalence
\[
    f^*_{\operatorname{PSh}}h_S(X)
    \simeq
    h_T(X_T).
\]
For every \(F\in\operatorname{PSh}(T)\) and
\(X\in\mathcal C_S\),
\[
    \bigl(
        f_{*,\operatorname{PSh}}F
    \bigr)(X)
    =
    F(X_T).
\]
\end{proposition}

\begin{proof}
Left Kan extension along a functor carries a representable presheaf to
the representable presheaf of its image.  The second formula is the
definition of restriction along \(u_f^{\mathrm{op}}\).
\end{proof}

\subsection{Descent to motivic spaces}

\begin{proposition}[Base change preserves motivic generators]
\label{prop:base-change-preserves-motivic-generators}
The functor
\[
    f^*_{\operatorname{PSh}}
\]
carries every morphism in \(W_{\mathrm{mot},S}\) to a motivic
equivalence in \(\operatorname{PSh}(T)\).
\end{proposition}

\begin{proof}
Let
\[
\begin{tikzcd}[column sep=large, row sep=large]
    W
        \arrow[r]
        \arrow[d]
    &
    V
        \arrow[d]
    \\
    U
        \arrow[r]
    &
    X
\end{tikzcd}
\]
be a chosen representative \(Q\) of a spectral Nisnevich distinguished
square over \(S\).

Since
\[
    f^*_{\operatorname{PSh}}
\]
preserves colimits and sends representables to their base changes, the
morphism
\[
    f^*_{\operatorname{PSh}}(e_Q)
\]
identifies with the excision morphism associated with the base-changed
square
\[
\begin{tikzcd}[column sep=large, row sep=large]
    W_T
        \arrow[r]
        \arrow[d]
    &
    V_T
        \arrow[d]
    \\
    U_T
        \arrow[r]
    &
    X_T.
\end{tikzcd}
\]
This square is spectral Nisnevich distinguished by
\Ref{prop:nisnevich-square-base-change}.

Its excision morphism is equivalent in the arrow category of
\(\operatorname{PSh}(T)\) to the excision morphism associated with the
chosen representative of its equivalence class. The latter belongs to
\(E_{\mathrm{Nis},T}\). Hence
\[
    f^*_{\operatorname{PSh}}(e_Q)
\]
is a motivic equivalence.

The initial presheaf is preserved by the left adjoint
\(f^*_{\operatorname{PSh}}\), and
\[
    f^*_{\operatorname{PSh}}h_S(\varnothing)
    \simeq
    h_T(\varnothing).
\]
Therefore
\[
    f^*_{\operatorname{PSh}}(e_{\varnothing})
\]
identifies with \(e_{\varnothing}\) over \(T\).

Finally, let
\[
    h_S(X\times_S\mathbb A^1_S)
    \longrightarrow
    h_S(X)
\]
be an affine-line projection. Using the formula for inverse image on
representables and the affine-space base-change equivalence, its image
is
\[
\begin{aligned}
    h_T
    \bigl(
        (X\times_S\mathbb A^1_S)\times_ST
    \bigr)
    &\longrightarrow
    h_T(X_T)
    \\
    \simeq\quad
    h_T
    \bigl(
        X_T\times_T\mathbb A^1_T
    \bigr)
    &\longrightarrow
    h_T(X_T).
\end{aligned}
\]
This belongs to \(W_{\mathbb A^1,T}\). Thus every morphism in
\(W_{\mathrm{mot},S}\) is carried to a motivic equivalence over \(T\).
\end{proof}

\begin{proposition}[Direct image preserves motivic-local objects]
\label{prop:direct-image-preserves-motivic-local}
If
\[
    F
    \in
    \mathbf H(T),
\]
then
\[
    f_{*,\operatorname{PSh}}F
\]
is motivic-local over \(S\).
\end{proposition}

\begin{proof}
Let \(Q\) be a spectral Nisnevich distinguished square in
\(\mathcal C_S\).  Evaluating \(f_{*,\operatorname{PSh}}F\) on \(Q\)
amounts to evaluating \(F\) on the base-changed square \(Q_T\), which is
distinguished by
\Ref{prop:nisnevich-square-base-change}.
Thus Nisnevich excision is preserved.  The empty value remains
contractible.

For \(X\in\mathcal C_S\), one has
\[
\begin{aligned}
    \bigl(
        f_{*,\operatorname{PSh}}F
    \bigr)
    (X\times_S\mathbb A^1_S)
    &=
    F
    \bigl(
        (X\times_S\mathbb A^1_S)\times_ST
    \bigr)
    \\
    &\simeq
    F
    \bigl(
        X_T\times_T\mathbb A^1_T
    \bigr)
    \\
    &\simeq
    F(X_T)
    \\
    &=
    \bigl(
        f_{*,\operatorname{PSh}}F
    \bigr)(X).
\end{aligned}
\]
Hence direct image preserves affine-line invariance.
\end{proof}

\begin{theorem}[Inverse and direct image]
\label{thm:motivic-inverse-direct-image}
The presheaf adjunction descends to an adjunction
\[
    f^*
    \colon
    \mathbf H(S)
    \rightleftarrows
    \mathbf H(T)
    :
    f_*.
\]
The inverse-image functor is colimit-preserving and satisfies
\[
    f^*\mathrm M_S(X)
    \simeq
    \mathrm M_T(X_T)
\]
for every \(X\in\mathcal C_S\).

The direct-image functor is given on sections by
\[
    (f_*F)(X)
    \simeq
    F(X_T)
\]
for every \(F\in\mathbf H(T)\) and \(X\in\mathcal C_S\).
\end{theorem}

\begin{proof}
Let
\[
    \mathcal V_f
\]
be the class of morphisms
\[
    u
    \colon
    F
    \longrightarrow
    G
\]
in \(\operatorname{PSh}(S)\) such that
\[
    f^*_{\operatorname{PSh}}(u)
\]
is a motivic equivalence in \(\operatorname{PSh}(T)\).

The class of motivic equivalences in
\(\operatorname{PSh}(T)\) is strongly saturated. Indeed, by
\Ref{thm:brave-new-motivic-localization}, it is the class of
morphisms inverted by the reflective localization
\[
    L_{\mathrm{mot},T}
    \colon
    \operatorname{PSh}(T)
    \longrightarrow
    \mathbf H(T).
\]

The functor
\[
    f^*_{\operatorname{PSh}}
    \colon
    \operatorname{PSh}(S)
    \longrightarrow
    \operatorname{PSh}(T)
\]
preserves colimits, compositions, and retracts. Consequently,
\(\mathcal V_f\) is strongly saturated.

By
\Ref{prop:base-change-preserves-motivic-generators},
the class \(\mathcal V_f\) contains every morphism in
\[
    W_{\mathrm{mot},S}.
\]
It therefore contains the strongly saturated class generated by
\(W_{\mathrm{mot},S}\). By
\Ref{thm:brave-new-motivic-localization}, this is precisely the class
of motivic equivalences in \(\operatorname{PSh}(S)\).

Thus
\[
    f^*_{\operatorname{PSh}}
\]
carries arbitrary motivic equivalences over \(S\) to motivic
equivalences over \(T\). The universal property of motivic localization
therefore supplies a unique colimit-preserving functor
\[
    f^*
    \colon
    \mathbf H(S)
    \longrightarrow
    \mathbf H(T)
\]
together with a canonical equivalence
\[
    f^*L_{\mathrm{mot},S}
    \simeq
    L_{\mathrm{mot},T}f^*_{\operatorname{PSh}}.
\]

For every \(X\in\mathcal C_S\), one has
\[
\begin{aligned}
    f^*\mathrm M_S(X)
    &=
    f^*L_{\mathrm{mot},S}h_S(X)
    \\
    &\simeq
    L_{\mathrm{mot},T}
    f^*_{\operatorname{PSh}}h_S(X)
    \\
    &\simeq
    L_{\mathrm{mot},T}h_T(X_T)
    \\
    &=
    \mathrm M_T(X_T),
\end{aligned}
\]
where the second equivalence uses
\Ref{prop:presheaf-base-change-formulas}.

By
\Ref{prop:direct-image-preserves-motivic-local},
the presheaf direct image restricts to a functor
\[
    f_*
    \colon
    \mathbf H(T)
    \longrightarrow
    \mathbf H(S),
\]
characterized by
\[
    \iota_Sf_*G
    \simeq
    f_{*,\operatorname{PSh}}\iota_TG
\]
for every \(G\in\mathbf H(T)\).

Let
\[
    F
    \in
    \mathbf H(S),
    \qquad
    G
    \in
    \mathbf H(T).
\]
Using the localization adjunction, the presheaf adjunction, and the
fact that
\[
    f_{*,\operatorname{PSh}}\iota_TG
\]
is motivic-local, we obtain natural equivalences
\[
\begin{aligned}
    \operatorname{Map}_{\mathbf H(T)}
    \bigl(
        f^*F,
        G
    \bigr)
    &\simeq
    \operatorname{Map}_{\mathbf H(T)}
    \bigl(
        L_{\mathrm{mot},T}
        f^*_{\operatorname{PSh}}\iota_SF,
        G
    \bigr)
    \\
    &\simeq
    \operatorname{Map}_{\operatorname{PSh}(T)}
    \bigl(
        f^*_{\operatorname{PSh}}\iota_SF,
        \iota_TG
    \bigr)
    \\
    &\simeq
    \operatorname{Map}_{\operatorname{PSh}(S)}
    \bigl(
        \iota_SF,
        f_{*,\operatorname{PSh}}\iota_TG
    \bigr)
    \\
    &\simeq
    \operatorname{Map}_{\operatorname{PSh}(S)}
    \bigl(
        \iota_SF,
        \iota_Sf_*G
    \bigr)
    \\
    &\simeq
    \operatorname{Map}_{\mathbf H(S)}
    \bigl(
        F,
        f_*G
    \bigr).
\end{aligned}
\]
Thus
\[
    f^*
    \dashv
    f_*.
\]

Finally, for every \(F\in\mathbf H(T)\) and
\(X\in\mathcal C_S\), the definition of presheaf direct image gives
\[
\begin{aligned}
    (f_*F)(X)
    &=
    \bigl(
        f_{*,\operatorname{PSh}}\iota_TF
    \bigr)(X)
    \\
    &=
    F(X_T).
\end{aligned}
\]
This proves the section formula.
\end{proof}

\begin{proposition}[Coherent functoriality]
\label{prop:motivic-coherent-base-functoriality}
For composable morphisms of quasi-compact quasi-separated spectral
schemes
\[
    R
    \xrightarrow{g}
    T
    \xrightarrow{f}
    S,
\]
there are canonical coherent equivalences
\[
    (f\circ g)^*
    \simeq
    g^*f^*
\]
and
\[
    (f\circ g)_*
    \simeq
    f_*g_*.
\]
For the identity morphism, the corresponding functors are canonically
equivalent to the identity.
\end{proposition}

\begin{proof}
The geometric base-change functors satisfy coherent composition by
\Ref{prop:smooth-site-geometric-base-change}.  The left Kan extensions
therefore satisfy the corresponding coherent equivalences.  Passing to
motivic localization preserves them.  The direct-image equivalences are
their right-adjoint mates, or follow directly from restriction of
presheaves.
\end{proof}

\subsection{Monoidality of inverse image}

\begin{theorem}[Symmetric monoidal inverse image]
\label{thm:motivic-inverse-image-symmetric-monoidal}
For every morphism
\[
    f
    \colon
    T
    \longrightarrow
    S
\]
of quasi-compact quasi-separated spectral schemes, the functor
\[
    f^*
    \colon
    \mathbf H(S)
    \longrightarrow
    \mathbf H(T)
\]
is strong symmetric monoidal for the Cartesian structures. Thus there
are natural coherent equivalences
\[
    f^*(F\times_SG)
    \simeq
    f^*F\times_Tf^*G
\]
and
\[
    f^*(\mathbf 1_S)
    \simeq
    \mathbf 1_T.
\]
Consequently \(f_*\) is canonically lax symmetric monoidal.
\end{theorem}

\begin{proof}
Since the Cartesian symmetric monoidal structure is determined by
finite products, the projections
\[
    F\times_SG
    \longrightarrow
    F
\]
and
\[
    F\times_SG
    \longrightarrow
    G
\]
determine a canonical comparison morphism
\[
    \mu_{F,G}
    \colon
    f^*(F\times_SG)
    \longrightarrow
    f^*F\times_Tf^*G.
\]

The terminal object of \(\mathbf H(S)\) is
\[
    \mathbf 1_S
    \simeq
    \mathrm M_S(S).
\]
Therefore
\[
\begin{aligned}
    f^*(\mathbf 1_S)
    &\simeq
    f^*\mathrm M_S(S)
    \\
    &\simeq
    \mathrm M_T(S\times_ST)
    \\
    &\simeq
    \mathrm M_T(T)
    \\
    &=
    \mathbf 1_T.
\end{aligned}
\]
Thus the canonical comparison
\[
    f^*(\mathbf 1_S)
    \longrightarrow
    \mathbf 1_T
\]
is an equivalence.

Now let
\[
    X,Y
    \in
    \mathcal C_S.
\]
Using
\Ref{cor:products-of-smooth-motives}
and the formula for inverse image on smooth motives, we obtain
\[
\begin{aligned}
    f^*
    \bigl(
        \mathrm M_S(X)\times_S\mathrm M_S(Y)
    \bigr)
    &\simeq
    f^*\mathrm M_S(X\times_SY)
    \\
    &\simeq
    \mathrm M_T
    \bigl(
        (X\times_SY)\times_ST
    \bigr)
    \\
    &\simeq
    \mathrm M_T(X_T\times_TY_T)
    \\
    &\simeq
    \mathrm M_T(X_T)
    \times_T
    \mathrm M_T(Y_T)
    \\
    &\simeq
    f^*\mathrm M_S(X)
    \times_T
    f^*\mathrm M_S(Y).
\end{aligned}
\]
Under these identifications, the morphism
\[
    \mu_{\mathrm M_S(X),\mathrm M_S(Y)}
\]
is the canonical equivalence induced by the two projections. Hence it
is an equivalence.

Both bifunctors
\[
    (F,G)
    \longmapsto
    f^*(F\times_SG)
\]
and
\[
    (F,G)
    \longmapsto
    f^*F\times_Tf^*G
\]
preserve colimits separately in each variable: \(f^*\) preserves
colimits, and Cartesian products in the motivic categories preserve
colimits separately by
\Ref{thm:cartesian-monoidal-motivic-localization}.

The smooth motives generate \(\mathbf H(S)\) under colimits by
\Ref{prop:brave-new-motivic-compact-generation}. Fixing one variable
and extending from the smooth motives in the other, and then repeating
in the second variable, shows that
\[
    \mu_{F,G}
\]
is an equivalence for all \(F,G\in\mathbf H(S)\).

The Cartesian symmetric monoidal coherence maps are uniquely determined
by the finite-product comparison maps. The equivalences above therefore
equip \(f^*\) with a canonical coherent strong symmetric monoidal
structure.

A right adjoint to a strong symmetric monoidal functor is canonically
lax symmetric monoidal. Hence \(f_*\) is canonically lax symmetric
monoidal.
\end{proof}

\subsection{Pointed base change}

\begin{proposition}[Base change on pointed motivic spaces]
\label{prop:pointed-motivic-base-change}
The functor \(f^*\) induces a strong symmetric monoidal functor
\[
    f^*_\bullet
    \colon
    \mathbf H_\bullet(S)
    \longrightarrow
    \mathbf H_\bullet(T)
\]
for the smash products.  It admits a right adjoint
\[
    f_{*,\bullet}.
\]
For every \(X\in\mathcal C_S\),
\[
    f^*_\bullet\mathrm M_{S,+}(X)
    \simeq
    \mathrm M_{T,+}(X_T).
\]
\end{proposition}

\begin{proof}
The functor \(f^*\) preserves colimits and the terminal object.  It
therefore induces a colimit-preserving functor on pointed objects by
applying \(f^*\) to the underlying object and basepoint.  Strong
symmetric monoidality for smash products follows from strong Cartesian
monoidality and the pushout formula defining the smash product.

Since the pointed motivic categories are presentable, the adjoint
functor theorem supplies a right adjoint.  Equivalently, the basepoint
of \(f_{*,\bullet}(G)\) is the composite
\[
    *_S
    \longrightarrow
    f_*(*_T)
    \longrightarrow
    f_*U_T(G)
\]
induced by the adjunction unit and the given basepoint of \(G\).
The formula on pointed smooth motives follows from
\Ref{thm:motivic-inverse-direct-image}.
\end{proof}

\begin{proposition}[Coherent pointed base functoriality]
\label{prop:pointed-motivic-coherent-base-functoriality}
For composable morphisms of quasi-compact quasi-separated spectral
schemes
\[
    R
    \xrightarrow{g}
    T
    \xrightarrow{f}
    S,
\]
there are canonical coherent equivalences
\[
    (f\circ g)^*_\bullet
    \simeq
    g^*_\bullet f^*_\bullet
\]
and
\[
    (f\circ g)_{*,\bullet}
    \simeq
    f_{*,\bullet}g_{*,\bullet}.
\]
For an identity morphism, the corresponding pointed inverse- and
direct-image functors are canonically equivalent to the identity.
\end{proposition}

\begin{proof}
The pointed inverse-image functor
\[
    f^*_\bullet
    \colon
    \mathbf H_\bullet(S)
    \longrightarrow
    \mathbf H_\bullet(T)
\]
is induced from the functor
\[
    f^*
    \colon
    \mathbf H(S)
    \longrightarrow
    \mathbf H(T)
\]
together with the canonical equivalence
\[
    f^*(*_S)
    \simeq
    *_T.
\]
The coherent composition equivalence
\[
    (f\circ g)^*
    \simeq
    g^*f^*
\]
of
\Ref{prop:motivic-coherent-base-functoriality}
therefore induces a coherent equivalence on pointed objects
\[
    (f\circ g)^*_\bullet
    \simeq
    g^*_\bullet f^*_\bullet.
\]
The identity equivalence is inherited in the same way.

The direct-image equivalence is the right-adjoint mate of the
inverse-image equivalence. Equivalently, both
\[
    (f\circ g)_{*,\bullet}
    \qquad\text{and}\qquad
    f_{*,\bullet}g_{*,\bullet}
\]
are right adjoint to the canonically equivalent functors
\[
    (f\circ g)^*_\bullet
    \qquad\text{and}\qquad
    g^*_\bullet f^*_\bullet.
\]
Uniqueness of right adjoints therefore supplies a canonical
equivalence
\[
    (f\circ g)_{*,\bullet}
    \simeq
    f_{*,\bullet}g_{*,\bullet}.
\]
Compatibility with triple composition and identities follows from the
coherence of the unpointed equivalences and the coherence of taking
adjoint mates.
\end{proof}

\section{Smooth direct image}
\label{sec:brave-new-motivic-smooth-direct-image}

Let
\[
    p
    \colon
    X
    \longrightarrow
    S
\]
be smooth and of finite presentation.  Thus
\[
    X
    \in
    \mathcal C_S.
\]

Since \(S\) is quasi-compact and quasi-separated and \(p\) is of finite
presentation, the spectral scheme \(X\) is quasi-compact and
quasi-separated. Consequently, the motivic categories
\[
    \mathbf H(X)
    \qquad\text{and}\qquad
    \mathbf H_\bullet(X)
\]
are defined under the standing conventions of this chapter.

\subsection{The geometric adjunction on smooth sites}

The geometric base-change functor is
\[
    p^*_{\mathrm{geom}}
    \colon
    \mathcal C_S
    \longrightarrow
    \mathcal C_X,
    \qquad
    Y
    \longmapsto
    Y\times_SX.
\]
The smooth-composition functor of
\Ref{prop:smooth-composition-functor}
is
\[
    p^{\mathrm{geom}}_\sharp
    \colon
    \mathcal C_X
    \longrightarrow
    \mathcal C_S,
\]
which sends \(Y\to X\) to its composite \(Y\to S\).

\begin{proposition}[Geometric smooth adjunction]
\label{prop:geometric-smooth-adjunction}
There is a natural adjunction
\[
    p^{\mathrm{geom}}_\sharp
    \dashv
    p^*_{\mathrm{geom}}.
\]
\end{proposition}

\begin{proof}
For \(Y\to X\) smooth of finite presentation and
\(Z\to S\) smooth of finite presentation, the universal property of the
fibre product gives
\[
    \operatorname{Map}_S(Y,Z)
    \simeq
    \operatorname{Map}_X
    \bigl(
        Y,
        Z\times_SX
    \bigr).
\]
This is the required adjunction.
\end{proof}

\subsection{Smooth direct image on presheaves}

\begin{construction}[Presheaf smooth direct image]
\label{con:presheaf-smooth-direct-image}
Define
\[
    p^{\operatorname{PSh}}_\sharp
    :=
    \operatorname{Lan}_{(p^{\mathrm{geom}}_\sharp)^{\mathrm{op}}}
    \colon
    \operatorname{PSh}(X)
    \longrightarrow
    \operatorname{PSh}(S).
\]
\end{construction}

\begin{proposition}[Presheaf smooth adjunction]
\label{prop:presheaf-smooth-adjunction}
There is a natural adjunction
\[
    p^{\operatorname{PSh}}_\sharp
    \dashv
    p^*_{\operatorname{PSh}}.
\]
For every \(Y\in\mathcal C_X\),
\[
    p^{\operatorname{PSh}}_\sharp h_X(Y)
    \simeq
    h_S(Y),
\]
where on the right \(Y\) is regarded as a smooth spectral \(S\)-scheme
by composition with \(p\).
\end{proposition}

\begin{proof}
By
\Ref{prop:geometric-smooth-adjunction},
the opposite functor
\[
    (p^*_{\mathrm{geom}})^{\mathrm{op}}
\]
is left adjoint to
\[
    (p^{\mathrm{geom}}_\sharp)^{\mathrm{op}}.
\]
Left Kan extension along a left adjoint is restriction along its right
adjoint.  Consequently
\[
    p^*_{\operatorname{PSh}}
    =
    \operatorname{Lan}_{(p^*_{\mathrm{geom}})^{\mathrm{op}}}
\]
is canonically equivalent to restriction along
\((p^{\mathrm{geom}}_\sharp)^{\mathrm{op}}\).  Its left adjoint is
therefore the displayed left Kan extension
\(p^{\operatorname{PSh}}_\sharp\).

The formula on representables is the usual effect of left Kan
extension on representable presheaves.
\end{proof}

\begin{proposition}[Smooth inverse image preserves motivic-local objects]
\label{prop:smooth-presheaf-inverse-image-preserves-local}
Let
\[
    p
    \colon
    X
    \longrightarrow
    S
\]
be smooth and of finite presentation. If
\[
    G
    \in
    \mathbf H(S),
\]
then
\[
    p^*_{\operatorname{PSh}}\iota_SG
\]
is motivic-local over \(X\).

Consequently, there is a canonical equivalence
\[
    p^*_{\operatorname{PSh}}\iota_SG
    \simeq
    \iota_Xp^*G.
\]
\end{proposition}

\begin{proof}
By
\Ref{prop:presheaf-smooth-adjunction},
the functor
\[
    p^*_{\operatorname{PSh}}
\]
is canonically equivalent to restriction along
\[
    (p^{\mathrm{geom}}_\sharp)^{\mathrm{op}}.
\]
Thus, for every
\[
    Y
    \in
    \mathcal C_X,
\]
there is a natural equivalence
\[
    \bigl(
        p^*_{\operatorname{PSh}}\iota_SG
    \bigr)(Y)
    \simeq
    G(Y),
\]
where on the right \(Y\) is regarded as a smooth spectral \(S\)-scheme
by composition with \(p\).

Let
\[
\begin{tikzcd}[column sep=large, row sep=large]
    W
        \arrow[r]
        \arrow[d]
    &
    V
        \arrow[d]
    \\
    U
        \arrow[r]
    &
    Y
\end{tikzcd}
\]
be a spectral Nisnevich distinguished square in \(\mathcal C_X\).
Regarded by composition as a square in \(\mathcal C_S\), it remains a
spectral Nisnevich distinguished square: its open immersion, étale
morphism, complementary closed immersion, and equivalence over the
closed complement are unchanged. Since \(G\) is Nisnevich-local, the
square
\[
\begin{tikzcd}[column sep=large, row sep=large]
    G(Y)
        \arrow[r]
        \arrow[d]
    &
    G(U)
        \arrow[d]
    \\
    G(V)
        \arrow[r]
    &
    G(W)
\end{tikzcd}
\]
is Cartesian. Hence
\[
    p^*_{\operatorname{PSh}}\iota_SG
\]
satisfies Nisnevich excision. Its value at the empty spectral scheme is
contractible because the empty spectral scheme is unchanged under
composition with \(p\).

Finally, for every \(Y\in\mathcal C_X\), the affine-space base-change
equivalence gives
\[
    \mathbb A^1_X
    \simeq
    X\times_S\mathbb A^1_S
\]
and therefore
\[
    Y\times_X\mathbb A^1_X
    \simeq
    Y\times_S\mathbb A^1_S.
\]
It follows that
\[
\begin{aligned}
    \bigl(
        p^*_{\operatorname{PSh}}\iota_SG
    \bigr)
    (Y\times_X\mathbb A^1_X)
    &\simeq
    G(Y\times_S\mathbb A^1_S)
    \\
    &\simeq
    G(Y)
    \\
    &\simeq
    \bigl(
        p^*_{\operatorname{PSh}}\iota_SG
    \bigr)(Y).
\end{aligned}
\]
Thus
\[
    p^*_{\operatorname{PSh}}\iota_SG
\]
is \(\mathbb A^1\)-invariant and hence motivic-local.

By the definition of motivic inverse image,
\[
    p^*G
    \simeq
    L_{\mathrm{mot},X}
    p^*_{\operatorname{PSh}}\iota_SG.
\]
Since the presheaf on the right is already motivic-local, the
localization unit is an equivalence. Therefore
\[
    p^*_{\operatorname{PSh}}\iota_SG
    \simeq
    \iota_Xp^*G.
\]
\end{proof}

\begin{proposition}[Smooth direct image preserves motivic generators]
\label{prop:smooth-direct-image-preserves-generators}
The functor
\[
    p^{\operatorname{PSh}}_\sharp
\]
carries every morphism in \(W_{\mathrm{mot},X}\) to a motivic
equivalence over \(S\).
\end{proposition}

\begin{proof}
Let
\[
\begin{tikzcd}[column sep=large, row sep=large]
    W
        \arrow[r]
        \arrow[d]
    &
    V
        \arrow[d]
    \\
    U
        \arrow[r]
    &
    Y
\end{tikzcd}
\]
be a chosen representative of a spectral Nisnevich distinguished square
in \(\mathcal C_X\).

Regarded by composition with \(p\) as a square in \(\mathcal C_S\), it
remains spectral Nisnevich distinguished: its open immersion, étale
morphism, complementary closed immersion, and equivalence over the
closed complement are unchanged.

The functor
\[
    p^{\operatorname{PSh}}_\sharp
\]
preserves colimits and sends each representable \(h_X(Z)\) to the
representable \(h_S(Z)\), where \(Z\) on the right is regarded as a
smooth spectral \(S\)-scheme by composition with \(p\). Consequently,
it carries the excision morphism of the displayed square over \(X\) to
the excision morphism of the same square regarded over \(S\).

This excision morphism is equivalent in the arrow category of
\(\operatorname{PSh}(S)\) to the excision morphism associated with the
chosen representative of its equivalence class. The latter belongs to
\(E_{\mathrm{Nis},S}\). Hence the image is a motivic equivalence.

The functor \(p^{\operatorname{PSh}}_\sharp\) preserves the initial
presheaf and sends \(h_X(\varnothing)\) to \(h_S(\varnothing)\).
Therefore it carries the empty excision morphism over \(X\) to the empty
excision morphism over \(S\).

Finally, let \(Y\in\mathcal C_X\). Base change of brave new affine
space gives
\[
    \mathbb A^1_X
    \simeq
    X\times_S\mathbb A^1_S.
\]
Therefore
\[
\begin{aligned}
    Y\times_X\mathbb A^1_X
    &\simeq
    Y\times_X
    \bigl(
        X\times_S\mathbb A^1_S
    \bigr)
    \\
    &\simeq
    Y\times_S\mathbb A^1_S.
\end{aligned}
\]
The affine-line projection over \(X\),
\[
    h_X(Y\times_X\mathbb A^1_X)
    \longrightarrow
    h_X(Y),
\]
is consequently carried to
\[
    h_S(Y\times_S\mathbb A^1_S)
    \longrightarrow
    h_S(Y),
\]
which belongs to \(W_{\mathbb A^1,S}\). Thus
\(p^{\operatorname{PSh}}_\sharp\) carries every morphism in
\(W_{\mathrm{mot},X}\) to a motivic equivalence over \(S\).
\end{proof}

\begin{theorem}[Smooth direct image on motivic spaces]
\label{thm:smooth-motivic-direct-image}
The presheaf smooth direct image descends to a colimit-preserving
functor
\[
    p_\sharp
    \colon
    \mathbf H(X)
    \longrightarrow
    \mathbf H(S)
\]
which is left adjoint to inverse image:
\[
    p_\sharp
    \dashv
    p^*.
\]
For every \(Y\in\mathcal C_X\),
\[
    p_\sharp\mathrm M_X(Y)
    \simeq
    \mathrm M_S(Y).
\]
\end{theorem}

\begin{proof}
By
\Ref{prop:smooth-direct-image-preserves-generators},
the colimit-preserving functor
\[
    p^{\operatorname{PSh}}_\sharp
\]
carries every morphism in \(W_{\mathrm{mot},X}\) to a motivic
equivalence over \(S\).

The class of morphisms carried by
\(p^{\operatorname{PSh}}_\sharp\) to motivic equivalences is strongly
saturated: the functor preserves colimits, compositions, and retracts,
and the class of motivic equivalences is strongly saturated.
Consequently,
\[
    p^{\operatorname{PSh}}_\sharp
\]
carries every motivic equivalence over \(X\) to a motivic equivalence
over \(S\).

It therefore descends uniquely to a colimit-preserving functor
\[
    p_\sharp
    \colon
    \mathbf H(X)
    \longrightarrow
    \mathbf H(S)
\]
characterized by a canonical equivalence
\[
    p_\sharp L_{\mathrm{mot},X}
    \simeq
    L_{\mathrm{mot},S}
    p^{\operatorname{PSh}}_\sharp.
\]

Let
\[
    F
    \in
    \mathbf H(X),
    \qquad
    G
    \in
    \mathbf H(S).
\]
Using the localization adjunction, the presheaf smooth adjunction, and
\Ref{prop:smooth-presheaf-inverse-image-preserves-local},
we obtain natural equivalences
\[
\begin{aligned}
    \operatorname{Map}_{\mathbf H(S)}
    \bigl(
        p_\sharp F,
        G
    \bigr)
    &\simeq
    \operatorname{Map}_{\operatorname{PSh}(S)}
    \bigl(
        p^{\operatorname{PSh}}_\sharp\iota_XF,
        \iota_SG
    \bigr)
    \\
    &\simeq
    \operatorname{Map}_{\operatorname{PSh}(X)}
    \bigl(
        \iota_XF,
        p^*_{\operatorname{PSh}}\iota_SG
    \bigr)
    \\
    &\simeq
    \operatorname{Map}_{\operatorname{PSh}(X)}
    \bigl(
        \iota_XF,
        \iota_Xp^*G
    \bigr)
    \\
    &\simeq
    \operatorname{Map}_{\mathbf H(X)}
    \bigl(
        F,
        p^*G
    \bigr).
\end{aligned}
\]
Thus
\[
    p_\sharp
    \dashv
    p^*.
\]

For \(Y\in\mathcal C_X\), the formula on representables from
\Ref{prop:presheaf-smooth-adjunction}
gives
\[
    p^{\operatorname{PSh}}_\sharp h_X(Y)
    \simeq
    h_S(Y),
\]
where \(Y\) on the right is regarded as a smooth spectral \(S\)-scheme
by composition with \(p\). Hence
\[
\begin{aligned}
    p_\sharp\mathrm M_X(Y)
    &\simeq
    p_\sharp L_{\mathrm{mot},X}h_X(Y)
    \\
    &\simeq
    L_{\mathrm{mot},S}
    p^{\operatorname{PSh}}_\sharp h_X(Y)
    \\
    &\simeq
    L_{\mathrm{mot},S}h_S(Y)
    \\
    &=
    \mathrm M_S(Y).
\end{aligned}
\]
\end{proof}

\begin{proposition}[Composition of smooth direct image]
\label{prop:smooth-direct-image-composition}
For smooth morphisms of finite presentation
\[
    Y
    \xrightarrow{q}
    X
    \xrightarrow{p}
    S,
\]
there is a canonical coherent equivalence
\[
    (p\circ q)_\sharp
    \simeq
    p_\sharp q_\sharp.
\]
The smooth direct image of an identity morphism is canonically the
identity.
\end{proposition}

\begin{proof}
At the geometric level, both functors are composition of smooth
spectral morphisms.  Their left Kan extensions therefore agree
canonically, with coherent associativity.  Passing to motivic
localization preserves this equivalence.
\end{proof}

\subsection{Smooth base change}

\begin{construction}[Smooth exchange transformation]
\label{con:unstable-smooth-exchange-transformation}
Consider a Cartesian square
\[
\begin{tikzcd}[column sep=large, row sep=large]
    X_T
        \arrow[r, "g'"]
        \arrow[d, "p'"']
    &
    X
        \arrow[d, "p"]
    \\
    T
        \arrow[r, "g"']
    &
    S
\end{tikzcd}
\]
in which \(p\) is smooth and of finite presentation.

Coherent functoriality of inverse image gives a canonical equivalence
\[
    (g')^*p^*
    \xrightarrow{\simeq}
    (p')^*g^*.
\]
Its left mate under the adjunctions
\[
    p_\sharp
    \dashv
    p^*
    \qquad\text{and}\qquad
    p'_\sharp
    \dashv
    (p')^*
\]
is the \textbf{smooth exchange transformation}
\[
    \operatorname{Ex}_{\sharp}^{*}
    \colon
    p'_\sharp(g')^*
    \longrightarrow
    g^*p_\sharp.
\]

Explicitly, for
\[
    F
    \in
    \mathbf H(X),
\]
the component
\[
    \operatorname{Ex}_{\sharp,F}^{*}
    \colon
    p'_\sharp(g')^*F
    \longrightarrow
    g^*p_\sharp F
\]
is adjoint to the composite
\[
\begin{aligned}
    (g')^*F
    &\xrightarrow{(g')^*(\eta_F)}
    (g')^*p^*p_\sharp F
    \\
    &\xrightarrow{\simeq}
    (p')^*g^*p_\sharp F,
\end{aligned}
\]
where
\[
    \eta_F
    \colon
    F
    \longrightarrow
    p^*p_\sharp F
\]
is the unit of the adjunction
\[
    p_\sharp
    \dashv
    p^*.
\]
\end{construction}

\begin{theorem}[Smooth base change]
\label{thm:unstable-smooth-base-change}
Let
\[
\begin{tikzcd}[column sep=large, row sep=large]
    X_T
        \arrow[r, "g'"]
        \arrow[d, "p'"']
    &
    X
        \arrow[d, "p"]
    \\
    T
        \arrow[r, "g"']
    &
    S
\end{tikzcd}
\]
be a Cartesian square of quasi-compact quasi-separated spectral schemes,
where \(p\) is smooth and of finite presentation. Then \(p'\) is smooth
and of finite presentation, and the smooth exchange transformation of
\Ref{con:unstable-smooth-exchange-transformation}
\[
    \operatorname{Ex}_{\sharp}^{*}
    \colon
    p'_\sharp(g')^*
    \longrightarrow
    g^*p_\sharp
\]
is an equivalence.
\end{theorem}

\begin{proof}
Smoothness and finite presentation of
\[
    p'
    \colon
    X_T
    \longrightarrow
    T
\]
follow from the stability of smooth morphisms of finite presentation
under base change.

Both functors
\[
    p'_\sharp(g')^*
    \qquad\text{and}\qquad
    g^*p_\sharp
\]
preserve colimits. The smooth motives
\[
    \mathrm M_X(Y),
    \qquad
    Y\in\mathcal C_X,
\]
generate \(\mathbf H(X)\) under colimits by
\Ref{prop:brave-new-motivic-compact-generation}.
It is therefore enough to prove that the exchange transformation is an
equivalence on every smooth motive
\[
    \mathrm M_X(Y).
\]

For \(Y\in\mathcal C_X\), one has
\[
\begin{aligned}
    p'_\sharp(g')^*\mathrm M_X(Y)
    &\simeq
    p'_\sharp
    \mathrm M_{X_T}
    \bigl(
        Y\times_XX_T
    \bigr)
    \\
    &\simeq
    \mathrm M_T
    \bigl(
        Y\times_XX_T
    \bigr).
\end{aligned}
\]
Since
\[
    X_T
    \simeq
    X\times_ST,
\]
there is a canonical equivalence
\[
    Y\times_XX_T
    \simeq
    Y\times_ST.
\]
Consequently,
\[
    p'_\sharp(g')^*\mathrm M_X(Y)
    \simeq
    \mathrm M_T(Y\times_ST).
\]

On the other hand,
\[
\begin{aligned}
    g^*p_\sharp\mathrm M_X(Y)
    &\simeq
    g^*\mathrm M_S(Y)
    \\
    &\simeq
    \mathrm M_T(Y\times_ST).
\end{aligned}
\]

We now identify the exchange component under these equivalences. By
\Ref{con:unstable-smooth-exchange-transformation}, its component at
\(\mathrm M_X(Y)\) is the mate of the composite
\[
\begin{aligned}
    (g')^*\mathrm M_X(Y)
    &\xrightarrow{(g')^*(\eta)}
    (g')^*p^*p_\sharp\mathrm M_X(Y)
    \\
    &\xrightarrow{\simeq}
    (p')^*g^*p_\sharp\mathrm M_X(Y),
\end{aligned}
\]
where
\[
    \eta
    \colon
    \mathrm M_X(Y)
    \longrightarrow
    p^*p_\sharp\mathrm M_X(Y)
\]
is the unit of
\[
    p_\sharp
    \dashv
    p^*.
\]

Equivalently, the exchange component is the composite
\[
\begin{aligned}
    p'_\sharp(g')^*\mathrm M_X(Y)
    &\longrightarrow
    p'_\sharp(p')^*g^*p_\sharp\mathrm M_X(Y)
    \\
    &\xrightarrow{\varepsilon}
    g^*p_\sharp\mathrm M_X(Y),
\end{aligned}
\]
where
\[
    \varepsilon
    \colon
    p'_\sharp(p')^*
    \longrightarrow
    \operatorname{id}_{\mathbf H(T)}
\]
is the counit of
\[
    p'_\sharp
    \dashv
    (p')^*.
\]

The unit
\[
    \mathrm M_X(Y)
    \longrightarrow
    p^*p_\sharp\mathrm M_X(Y)
    \simeq
    \mathrm M_X(Y\times_SX)
\]
is induced by the canonical \(X\)-morphism
\[
    Y
    \longrightarrow
    Y\times_SX
\]
whose components are the identity morphism of \(Y\) and the given
morphism \(Y\to X\).

After pullback along
\[
    g'
    \colon
    X_T
    \longrightarrow
    X,
\]
application of \(p'_\sharp\), and composition with the counit
\[
    p'_\sharp(p')^*
    \longrightarrow
    \operatorname{id},
\]
the resulting mate is induced by the canonical associativity
equivalence
\[
    Y\times_XX_T
    \xrightarrow{\simeq}
    Y\times_ST.
\]
Thus, under the identifications above, the component
\[
    \operatorname{Ex}_{\sharp,\mathrm M_X(Y)}^{*}
\]
is the canonical equivalence
\[
    \mathrm M_T
    \bigl(
        Y\times_XX_T
    \bigr)
    \xrightarrow{\simeq}
    \mathrm M_T(Y\times_ST).
\]

Thus the smooth exchange transformation is an equivalence on every
generating smooth motive. Since both sides preserve colimits, it is an
equivalence on every object of \(\mathbf H(X)\).
\end{proof}

\subsection{The smooth projection formula}

\begin{construction}[Projection transformation]
\label{con:unstable-smooth-projection-transformation}
Let
\[
    F\in\mathbf H(X),
    \qquad
    G\in\mathbf H(S).
\]
The unit of the adjunction
\[
    \eta_F
    \colon
    F
    \longrightarrow
    p^*p_\sharp F
\]
and the strong Cartesian symmetric monoidal structure on \(p^*\)
determine a morphism
\[
\begin{aligned}
    F\times_Xp^*G
    &\xrightarrow{\eta_F\times\id}
    p^*p_\sharp F\times_Xp^*G
    \\
    &\xrightarrow{\simeq}
    p^*
    \bigl(
        p_\sharp F\times_SG
    \bigr).
\end{aligned}
\]
Its adjoint under
\[
    p_\sharp
    \dashv
    p^*
\]
is the canonical projection transformation
\[
    p_\sharp
    \bigl(
        F\times_Xp^*G
    \bigr)
    \longrightarrow
    p_\sharp F\times_SG.
\]
\end{construction}

\begin{theorem}[Smooth projection formula]
\label{thm:unstable-smooth-projection-formula}
For every smooth morphism of finite presentation
\[
    p
    \colon
    X
    \longrightarrow
    S
\]
and all
\[
    F\in\mathbf H(X),
    \qquad
    G\in\mathbf H(S),
\]
the projection transformation is an equivalence:
\[
    p_\sharp
    \bigl(
        F\times_Xp^*G
    \bigr)
    \xrightarrow{\simeq}
    p_\sharp F\times_SG.
\]
\end{theorem}

\begin{proof}
Both sides preserve colimits separately in \(F\) and \(G\):
\(p_\sharp\) and \(p^*\) preserve colimits, and Cartesian products in
the motivic categories preserve colimits separately.  It is therefore
enough to take
\[
    F=\mathrm M_X(Y),
    \qquad
    G=\mathrm M_S(Z),
\]
with \(Y\in\mathcal C_X\) and \(Z\in\mathcal C_S\).

Using
\Ref{cor:products-of-smooth-motives},
one obtains
\[
\begin{aligned}
    \mathrm M_X(Y)
    \times_X
    p^*\mathrm M_S(Z)
    &\simeq
    \mathrm M_X
    \bigl(
        Y\times_X(Z\times_SX)
    \bigr)
    \\
    &\simeq
    \mathrm M_X(Y\times_SZ).
\end{aligned}
\]
Applying \(p_\sharp\) gives
\[
    \mathrm M_S(Y\times_SZ).
\]
On the other hand,
\[
\begin{aligned}
    p_\sharp\mathrm M_X(Y)
    \times_S
    \mathrm M_S(Z)
    &\simeq
    \mathrm M_S(Y)
    \times_S
    \mathrm M_S(Z)
    \\
    &\simeq
    \mathrm M_S(Y\times_SZ).
\end{aligned}
\]
The projection transformation identifies with the identity under these
equivalences.  The result follows.
\end{proof}

\subsection{Smooth direct image on pointed motivic spaces}

\begin{construction}[Pointed smooth direct image]
\label{con:pointed-smooth-direct-image}
Let
\[
    F
    =
    (*_X\to U_XF)
    \in
    \mathbf H_\bullet(X).
\]
Define
\[
    p_{\sharp,\bullet}F
    :=
    p_\sharp U_XF
    \amalg_{p_\sharp(*_X)}
    *_S.
\]
The map
\[
    p_\sharp(*_X)
    \longrightarrow
    p_\sharp U_XF
\]
is induced by the basepoint, and the map
\[
    p_\sharp(*_X)
    \longrightarrow
    *_S
\]
is the unique morphism to the terminal object.
\end{construction}

Since
\[
    *_X
    =
    \mathrm M_X(X),
\]
there is a canonical equivalence
\[
    p_\sharp(*_X)
    \simeq
    \mathrm M_S(X).
\]
Thus the construction may also be written
\[
    p_{\sharp,\bullet}F
    \simeq
    p_\sharp U_XF
    \amalg_{\mathrm M_S(X)}
    *_S.
\]

\begin{proposition}[Pointed smooth adjunction]
\label{prop:pointed-smooth-adjunction}
The functor
\[
    p_{\sharp,\bullet}
    \colon
    \mathbf H_\bullet(X)
    \longrightarrow
    \mathbf H_\bullet(S)
\]
is left adjoint to
\[
    p^*_\bullet
    \colon
    \mathbf H_\bullet(S)
    \longrightarrow
    \mathbf H_\bullet(X).
\]
For every \(Y\in\mathcal C_X\),
\[
    p_{\sharp,\bullet}\mathrm M_{X,+}(Y)
    \simeq
    \mathrm M_{S,+}(Y).
\]
\end{proposition}

\begin{proof}
The formula of
\Ref{con:pointed-smooth-direct-image}
is the standard left adjoint induced on pointed objects by an adjunction
\[
    p_\sharp
    \dashv
    p^*
\]
whose right adjoint preserves the terminal object.  Explicitly, a
pointed morphism
\[
    p_\sharp U_XF
    \amalg_{p_\sharp(*_X)}
    *_S
    \longrightarrow
    G
\]
is equivalent to a morphism
\[
    p_\sharp U_XF
    \longrightarrow
    U_SG
\]
whose restriction along \(p_\sharp(*_X)\) is the adjoint of the
basepoint of \(p^*_\bullet G\).  By the unpointed adjunction, this is
equivalent to a pointed morphism
\[
    F
    \longrightarrow
    p^*_\bullet G.
\]

For a free pointed motive,
\[
    \mathrm M_{X,+}(Y)
    =
    \mathrm M_X(Y)\amalg *_X.
\]
Since \(p_\sharp\) preserves coproducts,
\[
    p_\sharp U_X
    \bigl(
        \mathrm M_{X,+}(Y)
    \bigr)
    \simeq
    \mathrm M_S(Y)
    \amalg
    \mathrm M_S(X).
\]
Pushing out the second summand along
\(\mathrm M_S(X)\to *_S\) gives
\(\mathrm M_S(Y)\amalg *_S\), as required.
\end{proof}

\begin{proposition}[Composition of pointed smooth direct image]
\label{prop:pointed-smooth-direct-image-composition}
For smooth morphisms of finite presentation
\[
    Y
    \xrightarrow{q}
    X
    \xrightarrow{p}
    S,
\]
there is a canonical coherent equivalence
\[
    (p\circ q)_{\sharp,\bullet}
    \simeq
    p_{\sharp,\bullet}q_{\sharp,\bullet}.
\]
The pointed smooth direct image of an identity morphism is canonically
equivalent to the identity.
\end{proposition}

\begin{proof}
By
\Ref{prop:pointed-smooth-adjunction},
the functors
\[
    (p\circ q)_{\sharp,\bullet}
\]
and
\[
    p_{\sharp,\bullet}q_{\sharp,\bullet}
\]
are respectively left adjoint to
\[
    (p\circ q)^*_\bullet
\]
and
\[
    q^*_\bullet p^*_\bullet.
\]
By
\Ref{prop:pointed-motivic-coherent-base-functoriality},
there is a canonical coherent equivalence
\[
    (p\circ q)^*_\bullet
    \simeq
    q^*_\bullet p^*_\bullet.
\]
Uniqueness of left adjoints therefore supplies a canonical equivalence
\[
    (p\circ q)_{\sharp,\bullet}
    \simeq
    p_{\sharp,\bullet}q_{\sharp,\bullet}.
\]

The identity statement follows in the same way from the identity
equivalence for pointed inverse image. Compatibility with triple
composition and identities follows from the coherence of the
inverse-image equivalences and the coherence of taking adjoint mates.
\end{proof}

\begin{construction}[Pointed smooth exchange transformation]
\label{con:pointed-unstable-smooth-exchange-transformation}
Consider a Cartesian square
\[
\begin{tikzcd}[column sep=large, row sep=large]
    X_T
        \arrow[r, "g'"]
        \arrow[d, "p'"']
    &
    X
        \arrow[d, "p"]
    \\
    T
        \arrow[r, "g"']
    &
    S
\end{tikzcd}
\]
in which \(p\) is smooth and of finite presentation.

Coherent functoriality of pointed inverse image gives a canonical
equivalence
\[
    (g')^*_\bullet p^*_\bullet
    \xrightarrow{\simeq}
    (p')^*_\bullet g^*_\bullet.
\]
Its left mate under the pointed adjunctions
\[
    p_{\sharp,\bullet}
    \dashv
    p^*_\bullet
    \qquad\text{and}\qquad
    p'_{\sharp,\bullet}
    \dashv
    (p')^*_\bullet
\]
is the \textbf{pointed smooth exchange transformation}
\[
    \operatorname{Ex}_{\sharp,\bullet}^{*}
    \colon
    p'_{\sharp,\bullet}(g')^*_\bullet
    \longrightarrow
    g^*_\bullet p_{\sharp,\bullet}.
\]

Explicitly, for
\[
    F
    \in
    \mathbf H_\bullet(X),
\]
its component is adjoint to the composite
\[
\begin{aligned}
    (g')^*_\bullet F
    &\xrightarrow{(g')^*_\bullet(\eta_F)}
    (g')^*_\bullet
    p^*_\bullet
    p_{\sharp,\bullet}F
    \\
    &\xrightarrow{\simeq}
    (p')^*_\bullet
    g^*_\bullet
    p_{\sharp,\bullet}F,
\end{aligned}
\]
where
\[
    \eta_F
    \colon
    F
    \longrightarrow
    p^*_\bullet p_{\sharp,\bullet}F
\]
is the unit of the pointed adjunction.
\end{construction}

\begin{theorem}[Pointed smooth base change]
\label{thm:pointed-unstable-smooth-base-change}
In the Cartesian square of
\Ref{thm:unstable-smooth-base-change}, the pointed smooth exchange
transformation of
\Ref{con:pointed-unstable-smooth-exchange-transformation}
\[
    \operatorname{Ex}_{\sharp,\bullet}^{*}
    \colon
    p'_{\sharp,\bullet}(g')^*_\bullet
    \longrightarrow
    g^*_\bullet p_{\sharp,\bullet}
\]
is an equivalence.
\end{theorem}

\begin{proof}
Both functors
\[
    p'_{\sharp,\bullet}(g')^*_\bullet
    \qquad\text{and}\qquad
    g^*_\bullet p_{\sharp,\bullet}
\]
preserve colimits. By
\Ref{prop:pointed-motivic-generation}, the free pointed smooth motives
\[
    \mathrm M_{X,+}(Y),
    \qquad
    Y\in\mathcal C_X,
\]
generate \(\mathbf H_\bullet(X)\) under colimits. It is therefore enough
to evaluate the pointed exchange transformation on such a generator.

For \(Y\in\mathcal C_X\), pointed inverse image and pointed smooth
direct image give
\[
\begin{aligned}
    p'_{\sharp,\bullet}(g')^*_\bullet
    \mathrm M_{X,+}(Y)
    &\simeq
    p'_{\sharp,\bullet}
    \mathrm M_{X_T,+}
    \bigl(
        Y\times_XX_T
    \bigr)
    \\
    &\simeq
    \mathrm M_{T,+}
    \bigl(
        Y\times_XX_T
    \bigr)
    \\
    &\simeq
    \mathrm M_{T,+}(Y\times_ST).
\end{aligned}
\]

On the other hand,
\[
\begin{aligned}
    g^*_\bullet p_{\sharp,\bullet}
    \mathrm M_{X,+}(Y)
    &\simeq
    g^*_\bullet
    \mathrm M_{S,+}(Y)
    \\
    &\simeq
    \mathrm M_{T,+}(Y\times_ST).
\end{aligned}
\]

Under these identifications, the component of
\[
    \operatorname{Ex}_{\sharp,\bullet}^{*}
\]
is the canonical equivalence induced by
\[
    Y\times_XX_T
    \simeq
    Y\times_ST.
\]
This follows directly from its mate construction in
\Ref{con:pointed-unstable-smooth-exchange-transformation}; equivalently,
it is the free-pointed form of the unpointed exchange equivalence of
\Ref{thm:unstable-smooth-base-change}.

Thus the pointed smooth exchange transformation is an equivalence on
the generating free pointed smooth motives. Since both sides preserve
colimits, it is an equivalence on every object of
\(\mathbf H_\bullet(X)\).
\end{proof}

\begin{construction}[Pointed projection transformation]
\label{con:pointed-unstable-smooth-projection-transformation}
Let
\[
    F
    \in
    \mathbf H_\bullet(X),
    \qquad
    G
    \in
    \mathbf H_\bullet(S).
\]
The unit of the pointed adjunction
\[
    \eta_F
    \colon
    F
    \longrightarrow
    p^*_\bullet p_{\sharp,\bullet}F
\]
and the strong symmetric monoidal structure on \(p^*_\bullet\)
determine a morphism
\[
\begin{aligned}
    F\wedge_Xp^*_\bullet G
    &\xrightarrow{\eta_F\wedge\id}
    p^*_\bullet p_{\sharp,\bullet}F
    \wedge_X
    p^*_\bullet G
    \\
    &\xrightarrow{\simeq}
    p^*_\bullet
    \bigl(
        p_{\sharp,\bullet}F
        \wedge_S
        G
    \bigr).
\end{aligned}
\]
Its adjoint under
\[
    p_{\sharp,\bullet}
    \dashv
    p^*_\bullet
\]
is the pointed projection transformation
\[
    p_{\sharp,\bullet}
    \bigl(
        F\wedge_Xp^*_\bullet G
    \bigr)
    \longrightarrow
    p_{\sharp,\bullet}F
    \wedge_S
    G.
\]
\end{construction}

\begin{theorem}[Pointed smooth projection formula]
\label{thm:pointed-unstable-smooth-projection-formula}
For
\[
    F\in\mathbf H_\bullet(X),
    \qquad
    G\in\mathbf H_\bullet(S),
\]
the pointed projection transformation is an equivalence:
\[
    p_{\sharp,\bullet}
    \bigl(
        F\wedge_Xp^*_\bullet G
    \bigr)
    \xrightarrow{\simeq}
    p_{\sharp,\bullet}F
    \wedge_S
    G.
\]
\end{theorem}

\begin{proof}
Both sides preserve colimits separately in \(F\) and \(G\):
the functors
\[
    p_{\sharp,\bullet}
    \qquad\text{and}\qquad
    p^*_\bullet
\]
preserve colimits, and the smash products preserve colimits separately
in each variable.

By
\Ref{prop:pointed-motivic-generation},
the free pointed smooth motives generate the pointed motivic categories
under colimits. It is therefore enough to take
\[
    F
    =
    \mathrm M_{X,+}(Y),
    \qquad
    G
    =
    \mathrm M_{S,+}(Z),
\]
where
\[
    Y
    \in
    \mathcal C_X,
    \qquad
    Z
    \in
    \mathcal C_S.
\]

Using pointed base change and the smash-product formula for free pointed
smooth motives, we obtain
\[
\begin{aligned}
    \mathrm M_{X,+}(Y)
    \wedge_X
    p^*_\bullet\mathrm M_{S,+}(Z)
    &\simeq
    \mathrm M_{X,+}(Y)
    \wedge_X
    \mathrm M_{X,+}(Z\times_SX)
    \\
    &\simeq
    \mathrm M_{X,+}
    \bigl(
        Y\times_X(Z\times_SX)
    \bigr)
    \\
    &\simeq
    \mathrm M_{X,+}(Y\times_SZ).
\end{aligned}
\]
Applying \(p_{\sharp,\bullet}\) gives
\[
\begin{aligned}
    p_{\sharp,\bullet}
    \bigl(
        \mathrm M_{X,+}(Y)
        \wedge_X
        p^*_\bullet\mathrm M_{S,+}(Z)
    \bigr)
    &\simeq
    p_{\sharp,\bullet}
    \mathrm M_{X,+}(Y\times_SZ)
    \\
    &\simeq
    \mathrm M_{S,+}(Y\times_SZ).
\end{aligned}
\]

On the other hand,
\[
\begin{aligned}
    p_{\sharp,\bullet}\mathrm M_{X,+}(Y)
    \wedge_S
    \mathrm M_{S,+}(Z)
    &\simeq
    \mathrm M_{S,+}(Y)
    \wedge_S
    \mathrm M_{S,+}(Z)
    \\
    &\simeq
    \mathrm M_{S,+}(Y\times_SZ).
\end{aligned}
\]
Under these identifications, the pointed projection transformation of
\Ref{con:pointed-unstable-smooth-projection-transformation}
is the identity of
\[
    \mathrm M_{S,+}(Y\times_SZ).
\]
Hence it is an equivalence on the generating objects.

Since both sides preserve colimits separately in \(F\) and \(G\), the
pointed projection transformation is an equivalence for arbitrary
\(F\) and \(G\).
\end{proof}

\section{The unstable brave new motivic package}
\label{sec:unstable-brave-new-motivic-package}

\begin{theorem}[Brave new motivic spaces]
\label{thm:brave-new-motivic-spaces-package}
For every quasi-compact quasi-separated spectral scheme \(S\), there is
a compactly generated presentable \(\infty\)-category
\[
    \mathbf H(S)
\]
characterized as the full subcategory of
\[
    \operatorname{PSh}
    \bigl(
        \operatorname{Sm}^{\mathrm{fp}}_{/S}
    \bigr)
\]
spanned by the presheaves which satisfy spectral Nisnevich descent and
brave new affine-line invariance.

The following structures and properties are canonically available.
\begin{enumerate}[label=(\roman*)]
    \item The localization
    \[
        L_{\mathrm{mot},S}
        \colon
        \operatorname{PSh}(S)
        \longrightarrow
        \mathbf H(S)
    \]
    is accessible and Cartesian symmetric monoidal.

    \item The category \(\mathbf H(S)\) is Cartesian closed and is
    compactly generated by the motives
    \[
        \mathrm M_S(X),
        \qquad
        X\in\operatorname{Sm}^{\mathrm{fp}}_{/S}.
    \]

    \item The pointed category
    \[
        \mathbf H_\bullet(S)
    \]
    is presentably symmetric monoidal under the smash product and is
    generated under colimits by
    \[
        \mathrm M_{S,+}(X).
    \]

    \item Every morphism
    \[
        f
        \colon
        T
        \longrightarrow
        S
    \]
    of quasi-compact quasi-separated spectral schemes induces an
    adjunction
    \[
        f^*
        \colon
        \mathbf H(S)
        \rightleftarrows
        \mathbf H(T)
        :
        f_*,
    \]
    where \(f^*\) is colimit-preserving and strong Cartesian symmetric
    monoidal.

    \item If
    \[
        p
        \colon
        X
        \longrightarrow
        S
    \]
    is smooth and of finite presentation, then there is an adjunction
    \[
        p_\sharp
        \colon
        \mathbf H(X)
        \rightleftarrows
        \mathbf H(S)
        :
        p^*.
    \]

    \item Smooth base change and the smooth projection formula hold:
    for every Cartesian square
    \[
    \begin{tikzcd}[column sep=large, row sep=large]
        X_T
            \arrow[r, "g'"]
            \arrow[d, "p'"']
        &
        X
            \arrow[d, "p"]
        \\
        T
            \arrow[r, "g"']
        &
        S
    \end{tikzcd}
    \]
    with \(p\) smooth and of finite presentation, there are canonical
    equivalences
    \[
        p'_\sharp(g')^*
        \simeq
        g^*p_\sharp,
    \]
    and, for \(F\in\mathbf H(X)\) and \(G\in\mathbf H(S)\),
    \[
        p_\sharp
        \bigl(
            F\times_Xp^*G
        \bigr)
        \simeq
        p_\sharp F\times_SG.
    \]

    \item All preceding inverse-image, direct-image, smooth-direct-image,
    base-change, and projection-formula constructions admit pointed
    analogues, with Cartesian product replaced by smash product.
\end{enumerate}

The assignments
\[
    S
    \longmapsto
    \mathbf H(S)
\]
and
\[
    S
    \longmapsto
    \mathbf H_\bullet(S)
\]
are coherently functorial on quasi-compact quasi-separated spectral
schemes.  Thus brave new motivic spaces form an unstable
\[
    (*,\sharp,\times)
\]
formalism, and pointed brave new motivic spaces form its smash-monoidal
counterpart.
\end{theorem}

\begin{proof}
Existence and characterization of \(\mathbf H(S)\) are
\Ref{thm:brave-new-motivic-localization}.  Cartesian monoidality is
\Ref{thm:cartesian-monoidal-motivic-localization}, and compact
generation is
\Ref{prop:brave-new-motivic-compact-generation}.  The pointed
symmetric monoidal structure is
\Ref{thm:pointed-motivic-smash-monoidal}.

Inverse and direct image are
\Ref{thm:motivic-inverse-direct-image}, with coherent functoriality from
\Ref{prop:motivic-coherent-base-functoriality} and monoidality from
\Ref{thm:motivic-inverse-image-symmetric-monoidal}.  Pointed inverse and
direct image are
\Ref{prop:pointed-motivic-base-change}, with coherent functoriality from
\Ref{prop:pointed-motivic-coherent-base-functoriality}.

Smooth direct image is
\Ref{thm:smooth-motivic-direct-image}, with coherent composition from
\Ref{prop:smooth-direct-image-composition}.  Its pointed analogue is
\Ref{prop:pointed-smooth-adjunction}, with coherent composition from
\Ref{prop:pointed-smooth-direct-image-composition}.

Smooth base change and the projection formula are
\Ref{thm:unstable-smooth-base-change} and
\Ref{thm:unstable-smooth-projection-formula}.  Their pointed forms are
\Ref{thm:pointed-unstable-smooth-base-change} and
\Ref{thm:pointed-unstable-smooth-projection-formula}.

These equivalences are defined from the coherent geometric
base-change, composition, monoidal, and adjunction data by localization
and passage to adjoint mates. They therefore satisfy the required
composition and identity coherences.
\end{proof}

The constructions of the chapter may be summarized by
\[
    \operatorname{PSh}(S)
    =
    \operatorname{Fun}
    \left(
        \bigl(
            \operatorname{Sm}^{\mathrm{fp}}_{/S}
        \bigr)^{\mathrm{op}},
        \mathcal S
    \right),
\]
\[
    \mathbf H(S)
    =
    \operatorname{PSh}(S)
    \left[
        E_{\mathrm{Nis},S}^{-1},
        W_{\mathbb A^1,S}^{-1}
    \right],
\]
and
\[
    \mathbf H_\bullet(S)
    =
    \mathbf H(S)_{*/}.
\]
For arbitrary morphisms one has
\[
    f^*
    \dashv
    f_*,
\]
while for smooth morphisms of finite presentation one has
\[
    p_\sharp
    \dashv
    p^*.
\]
These adjunctions satisfy smooth base change and the smooth projection
formula, in both unpointed and pointed form.

The next chapter applies the lifting and tubular-neighbourhood theorems
of
\Ref{sec:smooth-etale-lifting-neighbourhoods}
to prove the Morel--Voevodsky localization theorem in spectral algebraic
geometry.  It then develops closed recollement, homotopy purity, and
derived nilpotent invariance.
